%% 4-2 角度と円、ラジアン（第13回：組 V4-2。ページ制限なし。偶数ページ）
%% ほかの巻が指している形：
%%   5-15「弧度法：角θ（rad）＝弧の長さ÷半径。半径1の円なら、角と弧の長さは同じ数。1周は2π rad（360°）。円の接線は、半径に垂直。」
%%   5-3「円の接線は、円に1点だけでふれる直線。接線は、ふれた点で半径に垂直。」
%%   5-15・5-17「弧の長さ＝半径×角（s＝rθ）」
%% この章で定義した部品（頭文字 bkivc）：
%%   \bkivcstep{番号}{見出し}  説明の段の見出し（丸数字＋太字の1行）
%%   \bkivcline{式}            式を、まん中に1行で置く
%%   \bkivcmidashi{見出し}     本文の中の小見出し（それだけで1行）
%%   \bkivcsector{(中心)}{半径}{はじめの角}{おわりの角}{色}  おうぎ形（面をぬって、ふちをかく）
%%   \bkivcright{(頂点)}{1本めの向きの角}{2本めの向きの角}  直角の印（2mm）
%%   \bkivcpolar{目もり mm}{円の数}  極座標の方眼（第14回 U-P）
\makeatletter
\newcommand\bkivcstep[2]{\par\addvspace{6pt}\Needspace{5\baselineskip}\noindent
  {\bkcircled[bkOrange]{#1}\enspace\sffamily\bfseries\color{bkOrange!75!black}#2}\par\nobreak\vspace{2pt}\noindent\ignorespaces}
\newcommand\bkivcline[1]{\par\addvspace{3pt}{\centering #1\par}\addvspace{3pt}\noindent\ignorespaces}
\newcommand\bkivcmidashi[1]{\par\addvspace{4pt}\Needspace{6\baselineskip}\noindent\textbf{#1}\par\noindent\ignorespaces}
\newcommand\bkivcsector[5]{\fill[#5!22] #1 -- ++(#3:#2) arc (#3:#4:#2) -- cycle;
  \draw[#5!85!black, line width=0.8pt] #1 -- ++(#3:#2) arc (#3:#4:#2) -- cycle;}
\newcommand\bkivcright[3]{\draw[bkInk, line width=0.5pt] ($#1+(#2:2.2)$) -- ($#1+(#2:2.2)+(#3:2.2)$) -- ($#1+(#3:2.2)$);}
% 極座標の方眼（第14回）：中心 (0,0)、1目もり #1 mm の円を #2 本と、30°ごとの線、角の数  \bkivcpolar{目もり}{円の数}
\newcommand\bkivcpolar[2]{%
  \foreach \k in {1,...,#2} {\draw[bkGray!55, line width=0.4pt] (0,0) circle ({\k*#1});}
  \foreach \a in {0,30,...,330} {\draw[bkGray!55, line width=0.4pt] (0,0) -- (\a:{#2*#1});
    \node[font=\sffamily\footnotesize, text=bkGray!70!black, inner sep=0pt] at (\a:{#2*#1+5.2}) {\a°};}
  \draw[bkInk, line width=0.8pt] (0,0) -- (0:{#2*#1});
  \fill[bkInk] (0,0) circle (0.8);}
\newcommand\bkivcpoint[1]{\par\addvspace{4pt}\noindent{\sffamily\bfseries\color{bkBlue!80!black}ポイント}\enspace #1\par\addvspace{4pt}}
\makeatother

\bkchapter{4}{2}{角度と円、ラジアン}{\bklevelrange{1}{2}}{角度と円、ラジアン}
\bkmichishirube{1-19,1-20,2-9,2-15,4-1}{4-6,4-8,4-15,5-3,5-15,5-17}

\begin{goal}
\begin{itemize}[label={\color{\bkTT}●}, leftmargin=1.6\zw, itemsep=0pt, topsep=1pt]
\item \bklevel{2}\enspace 弦・弧・おうぎ形がわかり、弧の長さとおうぎ形の面積を、中心角に比例させて求められる。
\item \bklevel{4}\enspace 角をラジアン（弧の長さ÷半径）で表し、度と換算して、$s=r\theta$ と $S=\dfrac{1}{2}r^2\theta$ を使える。
\item \bklevel{2}\enspace 接線が半径に垂直なわけを説明でき、DEG・RAD と極座標 $(r,\ \theta)$ を使える。
\end{itemize}
\end{goal}

\begin{hatena}
\noindent\begin{minipage}[c]{0.56\linewidth}
公園の回る遊具に、ミナさんとソラさんが乗っています。ミナさんは中心から\qty{1}{m}、ソラさんは中心から\qty{2}{m}の所です。遊具が、ちょうど $\dfrac{1}{4}$ 回転（90°）しました。

2人が回った角度は、どちらも90°です。でも、地面の上を進んだ道のりは、外側のソラさんのほうが長そうです。

では、「進んだ道のり」から「回った角度」を言いあてる方法は、あるでしょうか。角度を「長さ」ではかることはできるかな？
\end{minipage}\hfill
\begin{minipage}[c]{0.4\linewidth}\centering
\begin{tikzpicture}[x=1mm, y=1mm]
  % 回る遊具を真上から見た図。中心から 1 m（11 mm）と 2 m（22 mm）
  \fill[bkGray!12] (0,0) circle (25);
  \draw[bkGray, line width=0.7pt] (0,0) circle (25);
  \draw[bkGray, line width=0.5pt, dash pattern=on 2pt off 1.5pt] (0,0) -- (25,0);
  \draw[bkGray, line width=0.5pt, dash pattern=on 2pt off 1.5pt] (0,0) -- (0,25);
  \draw[bkOrange, line width=1.5pt, -{Stealth[length=3.5pt]}] (11,0) arc (0:86:11);
  \draw[bkBlue, line width=1.5pt, -{Stealth[length=3.5pt]}] (22,0) arc (0:87:22);
  \fill[bkOrange!80!black] (11,0) circle (1.4);
  \fill[bkBlue!80!black] (22,0) circle (1.4);
  \fill[bkInk] (0,0) circle (0.9);
  \draw[bkInk, line width=0.5pt] (4.5,0) arc (0:90:4.5);
  \node[font=\sffamily\footnotesize] at (45:7.6) {90°};
  \node[font=\sffamily\footnotesize, anchor=north] at (11,-1.8) {\qty{1}{m}};
  \node[font=\sffamily\footnotesize, anchor=north] at (22,-1.8) {\qty{2}{m}};
  \node[font=\sffamily\footnotesize, anchor=north, text=bkInk] at (0,-1.8) {中心};
  \node[font=\sffamily\footnotesize, text=bkOrange!80!black, anchor=south west] at (8.6,8.6) {ミナ};
  \node[font=\sffamily\footnotesize, text=bkBlue!80!black, anchor=south west] at (16.5,16.5) {ソラ};
\end{tikzpicture}
\end{minipage}
\end{hatena}

\noindent 角度は、分度器ではかって「何度」と言うのがふつうです（1-20）。この章では、角度を「弧の長さ÷半径」という長さのわり算で表す方法を作ります。この方法は、回るものの動きを計算するときに、とても便利です。第5巻で $\sin$ や $\cos$ を微分するときも、物理で回転を調べるときも、この方法を使います。

%% ---------------- やってみよう・思い出そう ----------------
\begin{yattemiyou}[ひもで、弧の長さをはかろう]
\noindent\begin{minipage}[c]{0.6\linewidth}
\textbf{用意するもの}\par
コンパス、分度器、定規、たこ糸（細いひも）、電卓
\begin{enumerate}[label={\bkcircled[bkGreen]{\arabic*}}, leftmargin=2.4\zw, itemsep=1pt, topsep=3pt]
\item コンパスで、同じ中心Oのまわりに、半径\qty{3}{cm}と半径\qty{5}{cm}の円をかこう。
\item 分度器を使って、Oの所に60°の角をかこう。角の2本の線は、外の円までのばそう。
\item 2本の線にはさまれた円の曲がった部分に、糸をそわせよう。はしに印をつけて、糸をまっすぐにのばし、定規で長さをはかろう。
\item 「はかった長さ÷半径」を、電卓で計算しよう。90°でもやってみよう。
\end{enumerate}
\end{minipage}\hfill
\begin{minipage}[c]{0.36\linewidth}\centering
\begin{tikzpicture}[x=1mm, y=1mm]
  % 半径 3 cm と 5 cm の円を、0.7 倍（21 mm と 35 mm）でかく。60°の部分だけ
  \draw[bkGray, line width=0.5pt] (0,0) -- (37,0);
  \draw[bkGray, line width=0.5pt] (0,0) -- (60:37);
  \draw[bkGray!70, line width=0.4pt] (-8:21) arc (-8:68:21);
  \draw[bkGray!70, line width=0.4pt] (-8:35) arc (-8:68:35);
  \draw[bkOrange, line width=1.6pt] (21,0) arc (0:60:21);
  \draw[bkBlue, line width=1.6pt] (35,0) arc (0:60:35);
  \draw[bkInk, line width=0.5pt] (6,0) arc (0:60:6);
  \node[font=\sffamily\footnotesize] at (30:9.2) {60°};
  \fill[bkInk] (0,0) circle (0.8);
  \node[font=\sffamily\footnotesize, anchor=east] at (-1,0) {O};
  \node[font=\sffamily\footnotesize, anchor=north] at (10.5,-1.4) {\qty{3}{cm}};
  \node[font=\sffamily\footnotesize, anchor=north] at (28,-1.4) {\qty{5}{cm}};
\end{tikzpicture}
\end{minipage}
\par\vspace{4pt}
ミナさんがはかった結果です。右のはしの列を、自分で計算してみよう。
\begin{bkcenter}
\begin{bktable}{colspec={Q[c,m,wd=18mm] Q[c,m,wd=20mm] Q[c,m,wd=30mm] Q[c,m,wd=38mm]}, rowsep=1.2pt, row{1}={font=\sffamily\small}}
角 & 半径 & 糸ではかった長さ & はかった長さ÷半径 \\
60° & \qty{3}{cm} & \qty{3.1}{cm} & \\
60° & \qty{5}{cm} & \qty{5.2}{cm} & \\
90° & \qty{3}{cm} & \qty{4.7}{cm} & \\
90° & \qty{5}{cm} & \qty{7.9}{cm} & \\
\end{bktable}
\end{bkcenter}
\end{yattemiyou}

\noindent やってみようの答え：60°では、$3.1\div3=1.03\cdots$、$5.2\div5=1.04$ です。90°では、$4.7\div3=1.56\cdots$、$7.9\div5=1.58$ です。半径がちがうので、はかった長さはちがいます。それなのに、同じ角なら「長さ÷半径」は、ほとんど同じ数になりました。小さなちがいは、糸ではかったときのずれです（2-6）。このふしぎを、この章の数学の道具③で解きあかします。

\begin{omoidasou}
\bkrecall{1-20}{角の大きさは、回った大きさ。1回転は360°、直角は90°。}
\bkrecall{2-9}{円周＝直径×$\pi$。$\pi=3.14159\cdots$ は、どんな大きさの円でも同じ数（円周率）。}
\bkrecall{2-15}{比例：$x$ が2倍、3倍になると、$y$ も2倍、3倍になる。式は $y=ax$。}
\bkrecall{2-10}{換算は「単位つきの1」をかける。$\qty{2}{min}\times\qty{60}{s/min}=\qty{120}{s}$（min が約分されて消える）。}
\bkrecall{4-1}{二等辺三角形の2つの底角は等しい。三角形の内角の和は180°。三角形の外角は、となりにない2つの内角の和。}
\end{omoidasou}

%% ---------------- 数学の道具 ----------------
\bksection{doguu}{数学の道具}

\bksubsection{1}{円のことば：弦・弧・おうぎ形}

\leavevmode\label{bk:4-2:kotoba}コンパスで円をかくとき、針をさした点が円の\textbf{中心}です。針からえんぴつの先までの長さが\textbf{半径}です（1-19）。円の上のどの点も、中心からの長さは、ちょうど半径です。円は「中心から同じ長さの点を、ぜんぶ集めた形」なのです。

円の中心を通って、円の上の2点を結んだ線分が\textbf{直径}です。直径は、半径2つ分の長さです。

\begin{bkfigure}
\begin{tikzpicture}[x=1mm, y=1mm]
  % 左：中心・半径・直径
  \begin{scope}[shift={(0,0)}]
    \draw[bkInk!70, line width=0.9pt] (0,0) circle (17);
    \draw[bkBlue, line width=1.3pt] (0,0) -- (40:17);
    \draw[bkOrange, line width=1.3pt] (-17,0) -- (17,0);
    \fill[bkInk] (0,0) circle (0.9);
    \node[font=\sffamily\footnotesize, anchor=north] at (0,-1.4) {中心O};
    \node[font=\sffamily\footnotesize, text=bkBlue!80!black, anchor=south east] at (40:9) {半径};
    \node[font=\sffamily\footnotesize, text=bkOrange!80!black, anchor=north] at (10,-1.4) {直径};
    \node[font=\sffamily\small\bfseries, anchor=north] at (0,-21) {中心・半径・直径};
  \end{scope}
  % まん中：弦と弧
  \begin{scope}[shift={(46,0)}]
    \draw[bkInk!70, line width=0.9pt] (0,0) circle (17);
    \draw[bkOrange, line width=2pt] (30:17) arc (30:150:17);
    \draw[bkBlue, line width=1.3pt] (150:17) -- (30:17);
    \fill[bkInk] (30:17) circle (0.9);
    \fill[bkInk] (150:17) circle (0.9);
    \fill[bkInk] (0,0) circle (0.9);
    \node[font=\sffamily\footnotesize, anchor=west] at ($(30:17)+(0.8,0)$) {B};
    \node[font=\sffamily\footnotesize, anchor=east] at ($(150:17)+(-0.8,0)$) {A};
    \node[font=\sffamily\footnotesize, text=bkBlue!80!black, anchor=north] at (0,6.6) {弦AB};
    \node[font=\sffamily\footnotesize, text=bkOrange!80!black, anchor=south] at (0,18.6) {弧AB};
    \node[font=\sffamily\small\bfseries, anchor=north] at (0,-21) {弦と弧};
  \end{scope}
  % 右：中心角とおうぎ形
  \begin{scope}[shift={(92,0)}]
    \draw[bkInk!70, line width=0.9pt] (0,0) circle (17);
    \bkivcsector{(0,0)}{17}{20}{95}{bkGreen}
    \fill[bkInk] (0,0) circle (0.9);
    \draw[bkInk, line width=0.6pt] (20:5) arc (20:95:5);
    \node[font=\sffamily\footnotesize, anchor=north] at (0,-1.4) {O};
    \node[font=\sffamily\footnotesize] at (62:9) {中心角};
    \draw[bkInk, line width=0.4pt] (40:13) -- (40:20.5);
    \node[font=\sffamily\footnotesize, anchor=south west, inner sep=1pt] at (40:20.5) {おうぎ形};
    \node[font=\sffamily\small\bfseries, anchor=north] at (0,-21) {中心角とおうぎ形};
  \end{scope}
\end{tikzpicture}
\bkfigcaption{円のことば。まん中の図のオレンジが弧、青が弦。右の図の緑の部分が、おうぎ形。}
\end{bkfigure}

\begin{itemize}[leftmargin=1.2\zw, itemsep=2pt, topsep=2pt, label={\color{\bkTT}●}]
\item 円の上の2点A、Bを結んだ線分を、\textbf{弦}（げん）といいます。弦ABとよびます。
\item 円の上のAからBまでの、曲がった部分を\textbf{弧}（こ）といいます。弧ABとよびます。
\item 2本の半径にはさまれた角を\textbf{中心角}といいます。
\item 2本の半径と弧でかこまれた形を\textbf{おうぎ形}といいます。ピザの1切れや、開いたせんすの形です。
\end{itemize}

\noindent\begin{minipage}[c]{0.66\linewidth}
「弧」と「弦」は、どちらも弓（ゆみ）から来たことばです。弓の、曲がった木の部分の形が弧です。弓にぴんとはった糸（つる）が弦です。弦は、弧の両はしをまっすぐに結んでいます。

A、Bの2点で円を分けると、短い弧と長い弧の2つができます。ことわりがなければ、弧ABは短いほうを指します。
\end{minipage}\hfill
\begin{minipage}[c]{0.3\linewidth}\centering
\begin{tikzpicture}[x=1mm, y=1mm]
  % 弓：木の部分（弧）と、つる（弦）
  \draw[bkOrange!80!black, line width=2.4pt, line cap=round] (60:20) arc (60:-60:20);
  \draw[bkBlue, line width=0.8pt] (60:20) -- (-60:20);
  \draw[bkInk!70, line width=0.8pt, -{Stealth[length=3pt]}] (2,0) -- (24,0);
  \node[font=\sffamily\footnotesize, text=bkOrange!80!black, anchor=west] at (21,0) {弧};
  \node[font=\sffamily\footnotesize, text=bkBlue!80!black, anchor=east] at (9,9) {弦};
\end{tikzpicture}
\end{minipage}

\bkivcmidashi{円周の長さは $2\pi r$}
円のまわりの長さ（円周）は、直径×$\pi$ でした（2-9）。$\pi$（パイ）は、どんな円でも同じ数 $3.14159\cdots$ です。直径は半径の2倍なので、半径を $r$ とすると、直径は $2\times r$ です。
\bkivcline{円周＝直径×$\pi$＝$2\times r\times\pi$＝$2\pi r$}
文字式では、数を先に、次に $\pi$、そのあとに文字を書くのが約束です。だから $2\pi r$ と書きます。半径\qty{5}{cm}の円なら、円周は $2\pi\times5=10\pi$ で、約\qty{31.4}{cm}です。

この本では、「約」のかわりに、記号 $\approx$ を使うことがあります。$\approx$ は「ほぼ等しい」という意味です（5-8でくわしく）。$10\pi\approx31.4$ のように書きます。

\begin{seikaku}
\textbf{この章の $r$ は半径}\par
半径を $r$ と書くのは、英語で半径を radius（レイディアス）というからです。3-15では、相関係数を $r$ と書きました。同じ文字でも、この章の $r$ は半径で、相関係数とは関係ありません。文字は、そのつど「何を表すか」を決めて使います。
\end{seikaku}

\bksubsection{2}{おうぎ形の弧と面積は、中心角に比例する}

\leavevmode\label{bk:4-2:hirei}半径\qty{6}{cm}の丸いピザを、中心から8等分します。1切れの中心角は、$360^{\circ}\div8=45^{\circ}$ です。1切れ、2切れ、3切れ……と、となりどうしをくっつけて、大きなおうぎ形を作ってみましょう。

\begin{bkfigure}
\begin{tikzpicture}[x=1mm, y=1mm]
  \foreach \n/\xx in {1/0, 2/34, 3/68, 4/104} {
    \begin{scope}[shift={(\xx,0)}]
      \draw[bkGray!60, line width=0.5pt, dash pattern=on 1.5pt off 1.5pt] (0,0) circle (14);
      \foreach \k in {1,...,\n} {\pgfmathsetmacro\aa{(\k-1)*45}\pgfmathsetmacro\bb{\k*45}\bkivcsector{(0,0)}{14}{\aa}{\bb}{bkOrange}}
      \fill[bkInk] (0,0) circle (0.7);
      \node[font=\sffamily\footnotesize, anchor=north] at (0,-16) {\n 切れ（$\the\numexpr\n*45\relax$°）};
    \end{scope}}
\end{tikzpicture}
\bkfigcaption{1切れの中心角は45°。切れの数を2倍、3倍にすると、中心角も2倍、3倍になる。}
\end{bkfigure}

8切れは、どれも同じ形・同じ大きさです。1切れを中心のまわりに回すと、となりの1切れにぴったり重なるからです（合同。4-1）。だから、どの1切れも、弧の長さは同じで、面積も同じです。

\begin{itemize}[leftmargin=1.2\zw, itemsep=2pt, topsep=2pt, label={\color{\bkTT}●}]
\item 2切れのおうぎ形は、中心角が2倍、弧の長さも2倍、面積も2倍です。
\item 3切れなら、中心角も弧も面積も、すべて3倍です。
\end{itemize}
中心角が2倍、3倍になると、弧の長さも面積も2倍、3倍になります。これは比例です（2-15）。円全体（8切れ）の円周は $2\pi\times6=12\pi\approx\qty{37.70}{cm}$ なので、1切れの弧は $12\pi\div8=1.5\pi\approx\qty{4.71}{cm}$ です。

\begin{bkcenter}
\begin{bktable}{colspec={Q[l,m,wd=30mm] *{5}{Q[c,m,wd=15mm]}}, rowsep=1.5pt}
切れの数 & 1 & 2 & 3 & 4 & 8 \\
中心角 & 45° & 90° & 135° & 180° & 360° \\
弧の長さ（cm） & 4.71 & 9.42 & 14.14 & 18.85 & 37.70 \\
面積（1切れの何倍） & 1 & 2 & 3 & 4 & 8 \\
\end{bktable}
\bktabcaption{半径\qty{6}{cm}のおうぎ形。弧の長さも面積も、中心角に比例する。}
\end{bkcenter}

\bkivcmidashi{度で書いた弧の長さの式}
円全体は、中心角360°のおうぎ形と考えられます。その弧の長さは、円周の $2\pi r$ です。中心角が $a$° のおうぎ形は、円全体の $\dfrac{a}{360}$ にあたります。比例なので、弧の長さも円周の $\dfrac{a}{360}$ 倍です。
\bkivcline{中心角 $a$° のおうぎ形の弧の長さ＝$2\pi r\times\dfrac{a}{360}$}
360でわるのは、1回転が360°だからです（1-20）。半径\qty{6}{cm}、中心角60°なら、$2\pi\times6\times\dfrac{60}{360}=12\pi\times\dfrac{1}{6}=2\pi$ で、約\qty{6.28}{cm}です。

面積も同じように、円全体の面積の $\dfrac{a}{360}$ 倍です。円の面積は $\pi r^2$（半径×半径×円周率。小学校で学んだ式です。なぜそうなるかは4-5でくわしく調べます）なので、中心角 $a$° のおうぎ形の面積は $\pi r^2\times\dfrac{a}{360}$ です。⑤では、ラジアンを使って、もっと簡単な式にします。

\begin{column}{角で地球の大きさをはかった人}
\noindent\begin{minipage}[c]{0.6\linewidth}
いまから2200年ほど前、エジプトにいたエラトステネスは、角と弧の長さの比例を使って、地球1周の長さを求めました。夏至の日の正午、南のシエネという町では、太陽が真上にあり、深い井戸の底まで日光がとどきました。同じ時刻、北のアレクサンドリアでは、地面に立てた棒のかげから、太陽の光が真上から7.2°かたむいてさしていることがわかりました。
\end{minipage}\hfill
\begin{minipage}[c]{0.36\linewidth}\centering
\begin{tikzpicture}[x=1mm, y=1mm]
  % 地球（半径 18）。シエネは真上（90°）、アレクサンドリアは 90°＋7.2°…図では見やすく 20°にしてかく
  \fill[bkBlue!10] (0,0) circle (18);
  \draw[bkBlue!70!black, line width=0.8pt] (0,0) circle (18);
  \draw[bkInk, line width=0.6pt] (0,0) -- (90:24);
  \draw[bkInk, line width=0.6pt] (0,0) -- (110:24);
  \foreach \xx in {-12,-4,4} {\draw[bkOrange, line width=0.7pt, -{Stealth[length=2.5pt]}] (\xx,34) -- (\xx,{sqrt(18*18-\xx*\xx)+0.6});}
  \draw[bkInk, line width=0.5pt] (90:6) arc (90:110:6);
  \node[font=\sffamily\footnotesize, anchor=south] at (100:6.6) {7.2°};
  \node[font=\sffamily\footnotesize, anchor=south west] at (1,24.2) {シエネ};
  \node[font=\sffamily\footnotesize, anchor=south, align=center] at (112:24.5) {アレクサン\\ドリア};
  \fill[bkInk] (0,0) circle (0.7);
\end{tikzpicture}
\end{minipage}

\par\vspace{3pt}
太陽の光は平行にとどくので、7.2°は、地球の中心から見た2つの町の間の中心角と同じです（4-1の平行線と角）。$360\div7.2=50$ なので、2つの町の間の道のりは、地球1周の $\dfrac{1}{50}$ です。その道のりを約\qty{925}{km}とすると、地球1周は $925\times50=\num{46250}$ で、約\qty{46000}{km}。本当の値の約\qty{40000}{km}に近い数を、棒と角度と道のりだけで求めたのです。（図は、角を大きめにかいています。）
\end{column}

\bksubsection{3}{角を長さではかる：弧度法とラジアン}

\leavevmode\label{bk:4-2:rad}\bkhatenaref にもどりましょう。遊具が90°回るあいだに、2人はそれぞれ、半径1\,mと2\,mの円の $\dfrac{90}{360}=\dfrac{1}{4}$ を進みます。

\begin{bkcenter}
\begin{bktable}{colspec={Q[l,m,wd=18mm] Q[c,m,wd=16mm] Q[c,m,wd=48mm] Q[c,m,wd=34mm]}, rowsep=1.5pt, row{1}={font=\sffamily\small}}
 & 半径 & 進んだ道のり（弧の長さ） & 弧の長さ÷半径 \\
ミナ & \qty{1}{m} & $2\pi\times1\times\dfrac{1}{4}\approx\qty{1.57}{m}$ & $1.57\div1=1.57$ \\
ソラ & \qty{2}{m} & $2\pi\times2\times\dfrac{1}{4}\approx\qty{3.14}{m}$ & $3.14\div2=1.57$ \\
\end{bktable}
\bktabcaption{進んだ道のりは2倍ちがうのに、「弧の長さ÷半径」は同じ。}
\end{bkcenter}

ソラさんは、ミナさんの2倍の道のりを進みました。ところが、「弧の長さ÷半径」を計算すると、2人とも1.57で同じです。やってみようでも、同じことが起きていました。

\bkivcstep{1}{なぜ同じになるの？（図で考える）}
\noindent\begin{minipage}[c]{0.5\linewidth}
同じ中心のまわりに、半径1、2、3の円をかきます。同じ角のおうぎ形を切りとると、大きいおうぎ形は、小さいおうぎ形を、そのまま2倍、3倍に引きのばした形になっています。

引きのばすと、弧の長さも2倍、3倍になります。半径が2倍なら弧も2倍、半径が3倍なら弧も3倍。だから「弧÷半径」は変わりません。
\end{minipage}\hfill
\begin{minipage}[c]{0.46\linewidth}\centering
\begin{tikzpicture}[x=1mm, y=1mm]
  % 半径 1・2・3（12・24・36 mm）。角 50°
  \draw[bkGray, line width=0.5pt] (0,0) -- (40,0);
  \draw[bkGray, line width=0.5pt] (0,0) -- (50:40);
  \draw[bkPurple, line width=1.6pt] (36,0) arc (0:50:36);
  \draw[bkBlue, line width=1.6pt] (24,0) arc (0:50:24);
  \draw[bkOrange, line width=1.6pt] (12,0) arc (0:50:12);
  \fill[bkInk] (0,0) circle (0.8);
  \node[font=\sffamily\footnotesize, anchor=north] at (12,-1.4) {1};
  \node[font=\sffamily\footnotesize, anchor=north] at (24,-1.4) {2};
  \node[font=\sffamily\footnotesize, anchor=north] at (36,-1.4) {3};
  \node[font=\sffamily\footnotesize, anchor=north] at (0,-1.4) {半径};
  \node[font=\sffamily\footnotesize, text=bkOrange!80!black, anchor=west] at (14:13.2) {1こ分};
  \node[font=\sffamily\footnotesize, text=bkBlue!80!black, anchor=west] at (14:25.2) {2こ分};
  \node[font=\sffamily\footnotesize, text=bkPurple!80!black, anchor=west] at (14:37.2) {3こ分};
\end{tikzpicture}
\end{minipage}

\bkivcstep{2}{なぜ同じになるの？（式で考える）}
中心角 $a$° の弧の長さは、②の式で $2\pi r\times\dfrac{a}{360}$ でした。これを半径 $r$ でわります。
\bkivcline{$\left(2\pi r\times\dfrac{a}{360}\right)\div r=2\pi\times\dfrac{a}{360}$}
$r$ をかけてから $r$ でわるので、$r$ は消えます。のこった $2\pi\times\dfrac{a}{360}$ には、半径 $r$ が入っていません。角 $a$ だけで決まる数です。だから、どんな大きさの円で計算しても、同じ角なら同じ数になるのです。

\bkivcstep{3}{角を「弧の長さ÷半径」で表す}
「弧の長さ÷半径」は、角だけで決まる数でした。そこで、この数を、角の大きさそのものとして使うことにします。
\bkivcline{\fbox{角 $\theta$（ラジアン）＝弧の長さ÷半径}}
$\theta$ は、ギリシャ文字で「シータ」と読みます。角の大きさを表すのによく使う文字です。このように角をはかる方法を\textbf{弧度法}（こどほう）といい、この角の単位を\textbf{ラジアン}（記号 rad）といいます。\bkidx{こどほう}{弧度法}\bkidx{らじあん}{ラジアン}ラジアン（radian）は、半径（radius）から作られたことばです。

\bkhatenaref の答え：角度は、「弧の長さ÷半径」という長さのわり算で表せます。ミナさんもソラさんも、$\dfrac{\pi}{2}\approx1.57$ ラジアン回りました。進んだ道のりがちがっても、回った角は同じ数になります。

\bkivcstep{4}{半径1の円なら、角と弧の長さは同じ数}
半径が1のときは、1でわっても数は変わりません。だから、
\bkivcline{\textbf{半径1の円なら、角（rad）と弧の長さは同じ数}}
です。半径1の円の上を、長さ0.5だけ進んだら、角は0.5ラジアン。長さ2だけ進んだら、角は2ラジアンです。ラジアンは、「半径1の円の上を、どれだけ歩いたか」で角をはかる方法、と考えることもできます。

\bkivcstep{5}{1周は $2\pi$ ラジアン}
1周（360°）の弧の長さは、円周の $2\pi r$ です。半径でわると、
\bkivcline{1周の角＝$2\pi r\div r=2\pi$（ラジアン）}
です。半分（180°）なら $\pi$、$\dfrac{1}{4}$ 周（90°）なら $\dfrac{\pi}{2}$ です。

\begin{bkfigure}
\begin{tikzpicture}[x=1mm, y=1mm]
  % 左：1\,rad（弧の長さ＝半径）
  \begin{scope}[shift={(0,0)}]
    \draw[bkInk!50, line width=0.6pt] (0,0) circle (20);
    \draw[bkBlue, line width=1.8pt] (0,0) -- (20,0);
    \draw[bkBlue, line width=1.8pt] (20,0) arc (0:57.2958:20);
    \draw[bkInk, line width=0.7pt] (0,0) -- (57.2958:20);
    \draw[bkInk, line width=0.5pt] (5,0) arc (0:57.2958:5);
    \node[font=\sffamily\footnotesize] at (28.6:8.2) {1 rad};
    \node[font=\sffamily\footnotesize, text=bkBlue!80!black, anchor=north] at (10,-1.4) {半径 $r$};
    \node[font=\sffamily\footnotesize, text=bkBlue!80!black, anchor=west] at (28.6:21) {弧も $r$};
    \fill[bkInk] (0,0) circle (0.8);
    \node[font=\sffamily\small\bfseries, anchor=north] at (0,-24) {1 rad（約57.3°）};
  \end{scope}
  % 右：半径6本分と、あと0.28
  \begin{scope}[shift={(66,0)}]
    \foreach \k in {0,...,5} {\pgfmathsetmacro\c{mod(\k,2)==0 ? 1 : 0}
      \ifdim\c pt>0.5pt \draw[bkBlue, line width=2.2pt] ({\k*57.2958}:20) arc ({\k*57.2958}:{(\k+1)*57.2958-1.2}:20);
      \else \draw[bkOrange, line width=2.2pt] ({\k*57.2958}:20) arc ({\k*57.2958}:{(\k+1)*57.2958-1.2}:20); \fi
      \node[font=\sffamily\footnotesize] at ({(\k+0.5)*57.2958}:15.4) {\the\numexpr\k+1\relax};}
    \draw[bkGray, line width=2.2pt] (343.775:20) arc (343.775:360:20);
    \draw[bkInk!50, line width=0.6pt] (0,0) -- (20,0);
    \fill[bkInk] (0,0) circle (0.8);
    \node[font=\sffamily\footnotesize, text=bkGray, anchor=west] at (352:21.2) {あと0.28};
    \node[font=\sffamily\small\bfseries, anchor=north] at (0,-24) {1周は $2\pi\approx6.28$ rad};
  \end{scope}
\end{tikzpicture}
\bkfigcaption{左：弧の長さが半径と同じになる角が1\,rad。右：円周には、半径と同じ長さの弧が6こと、あと0.28こ分入る。}
\end{bkfigure}

\noindent\textbf{1ラジアンは、どのくらいの角？}\par
1ラジアンは、「弧の長さ÷半径」がちょうど1になる角です。つまり、\textbf{弧の長さが半径と同じになる角}です。半径と同じ長さのひもを、円にそわせて曲げると、その両はしと中心が作る角が1ラジアンです（図の左）。

円周は $2\pi r$ なので、半径の $2\pi\approx6.28$ 倍です。円周には、半径と同じ長さの弧が6こと、あと少し（0.28こ分）入ります（図の右）。1周360°を6.28こに分けるので、1ラジアンは $360\div6.28\approx57.3$ で、約57.3°です。1ラジアンは、1°とくらべると、ずっと大きな角です。

\bksubsection{4}{度とラジアンの換算}

\leavevmode\label{bk:4-2:kansan}度とラジアンは、同じ角を、ちがう単位で表したものです。cmとmのように、換算できます（2-10）。もとになるのは、半周の角です。
\bkivcline{\fbox{$\pi$ ラジアン＝180°}}

\bkivcmidashi{1°は何ラジアン？　1ラジアンは何度？}
両辺を180でわると、1°が何ラジアンかがわかります。両辺を $\pi$ でわると、1ラジアンが何度かがわかります。
\bkivcline{$1^{\circ}=\dfrac{\pi}{180}\ \text{rad}\approx\qty{0.01745}{rad}$\qquad $\qty{1}{rad}=\dfrac{180}{\pi}^{\circ}\approx57.30^{\circ}$}

\bkivcmidashi{換算は「単位つきの1」をかける}
2-10では、\qty{60}{s/min}のような「単位つきの1」をかけて、単位を換算しました。$\pi$ ラジアンと180°は同じ角なので、次の2つの分数は、どちらも1です。
\bkivcline{$\dfrac{\pi\ \text{rad}}{180^{\circ}}=1$\qquad\qquad $\dfrac{180^{\circ}}{\pi\ \text{rad}}=1$}
度をラジアンにするときは、°が約分されて消えるように、左の分数をかけます。ラジアンを度にするときは、rad が消えるように、右の分数をかけます。
\begin{align*}
60^{\circ}&=60^{\circ}\times\frac{\pi\ \text{rad}}{180^{\circ}}=\frac{60\pi}{180}\ \text{rad}=\frac{\pi}{3}\ \text{rad} &&\text{\small（60と180を60で約分）}\\
\frac{\pi}{4}\ \text{rad}&=\frac{\pi}{4}\ \text{rad}\times\frac{180^{\circ}}{\pi\ \text{rad}}=\frac{180^{\circ}}{4}=45^{\circ} &&\text{\small（$\pi$ が約分されて消える）}
\end{align*}
どちらの分数をかけるか迷ったら、「消したい単位が分母にくるほう」と考えます。

\begin{bkcenter}
\begin{bktable}{colspec={Q[l,m,wd=15mm] *{9}{Q[c,m,wd=11mm]}}, rowsep=1.5pt, colsep=1.5pt}
度 & 0° & 30° & 45° & 60° & 90° & 120° & 180° & 270° & 360° \\
ラジアン & 0 & $\dfrac{\pi}{6}$ & $\dfrac{\pi}{4}$ & $\dfrac{\pi}{3}$ & $\dfrac{\pi}{2}$ & $\dfrac{2\pi}{3}$ & $\pi$ & $\dfrac{3\pi}{2}$ & $2\pi$ \\
小数で & 0 & 0.524 & 0.785 & 1.047 & 1.571 & 2.094 & 3.142 & 4.712 & 6.283 \\
\end{bktable}
\bktabcaption{よく使う角。180°が $\pi$ なので、90°はその半分の $\dfrac{\pi}{2}$、60°は3分の1の $\dfrac{\pi}{3}$。}
\end{bkcenter}

表を全部おぼえる必要はありません。「180°が $\pi$」だけおぼえておけば、30°は180°の $\dfrac{1}{6}$ だから $\dfrac{\pi}{6}$、のように、いつでも作れます。

\begin{yattemiyou}[ラジアンの分度器を作ろう]
\noindent\begin{minipage}[c]{0.6\linewidth}
\textbf{用意するもの}\par
厚紙、コンパス、たこ糸（または紙テープ）、定規、はさみ、ふつうの分度器
\begin{enumerate}[label={\bkcircled[bkGreen]{\arabic*}}, leftmargin=2.4\zw, itemsep=1pt, topsep=3pt]
\item 厚紙に、半径\qty{10}{cm}の円をかこう。
\item 糸を、半径と同じ\qty{10}{cm}に切ろう。
\item 円の右はしから、糸を円にそわせて置き、糸のはしに印をつけて「1」と書こう。そこからまた糸をそわせて、「2」「3」……と印をつけよう。
\item 中心と印を結ぶ線を引くと、1\,rad、2\,rad、3\,rad……の目もりができる。6の印のあと、1周までに少しだけあまることを確かめよう。
\item ふつうの分度器を重ねて、1\,rad が何度か読もう。
\end{enumerate}
\end{minipage}\hfill
\begin{minipage}[c]{0.36\linewidth}\centering
\begin{tikzpicture}[x=0.85mm, y=0.85mm]
  \fill[bkYellow!15] (0,0) circle (20);
  \draw[bkInk!60, line width=0.8pt] (0,0) circle (20);
  \foreach \k in {0,...,6} {\draw[bkInk!70, line width=0.5pt] (0,0) -- ({\k*57.2958}:20);
    \draw[bkRed, line width=1pt] ({\k*57.2958}:18.5) -- ({\k*57.2958}:21.5);
    \node[font=\sffamily\footnotesize] at ({\k*57.2958}:24.5) {\k};}
  \draw[bkOrange, line width=1.6pt] (20,0) arc (0:57.3:20);
  \fill[bkInk] (0,0) circle (0.8);
\end{tikzpicture}
\end{minipage}
\par\vspace{3pt}
⑤の答え：1\,rad は、分度器で約57°です。どの大きさの円で作っても、糸を半径と同じ長さにすれば、同じ目もりになります。
\end{yattemiyou}

\begin{seikaku}
\textbf{ラジアンには、単位をつけないことが多い}\par
ラジアンは「長さ÷長さ」です。m÷mのように同じ単位どうしでわると、単位は約分されて消えます。だから、ラジアンは本当は単位のない数です。そのため、「$\theta=2$」と書いたら、ふつうは2ラジアンのことです。いっぽう、度は「1回転を360に分けたもの」という約束で決めた単位なので、「°」を省いてはいけません。
\end{seikaku}

\bkivcmidashi{わざわざラジアンを使うのは、なぜ？}
度は、1回転を360に分けた単位です。360という数は、昔の人が「2、3、4、5、6、8、9、10、12……と、たくさんの数でわり切れて便利だから」と選んだ数です。円の形とは関係がありません。だから、度で弧の長さを計算すると、$2\pi r\times\dfrac{a}{360}$ のように、360や $2\pi$ が入ってきます。ラジアンは、円の半径と弧の長さから決めた角です。次の⑤で見るように、ラジアンを使うと、式から360も $2\pi$ も消えて、とても簡単になります。

\begin{motto}[2]
\raggedright
\textbf{角の細かい単位「分」「秒」と、海の上の距離}\par
1°より小さい角は、1°を60に分けた\textbf{分}（′）と、1分を60に分けた\textbf{秒}（″）で表すことがあります。時計の分・秒と同じ60ずつの区切りです。地図の緯度（いど）や経度（けいど）は、北緯34度41分のように書きます。

船や飛行機では、地球の中心から見て緯度1分の角に当たる地面の弧の長さを、1海里（かいり）とよび、距離の単位にしています。地球の1周（約\qty{40000}{km}）を、$360\times60=\num{21600}$ 分でわると、$\num{40000000}\div\num{21600}\approx1852$ で、約\qty{1852}{m}です。いまは、1海里＝\qty{1852}{m}と決めています。角と弧の長さが比例する（②）ことを、そのまま使った単位です。
\end{motto}

\bksubsection{5}{弧の長さ $s=r\theta$ と、おうぎ形の面積}

\leavevmode\label{bk:4-2:srt}ラジアンの決め方「$\theta$＝弧の長さ÷半径」で、弧の長さを $s$ と書くと、$\theta=\dfrac{s}{r}$ です。$s$ は、弧の長さによく使う文字です。両辺に $r$ をかけると（2-14）、
\bkivcline{\fbox{弧の長さ＝半径×角\qquad $s=r\theta$\qquad（$\theta$ はラジアン）}}
になります。半径\qty{6}{cm}、中心角 $\dfrac{\pi}{3}$ ラジアン（60°）なら、弧の長さは
\bkivcline{$s=6\times\dfrac{\pi}{3}=2\pi\approx\qty{6.28}{cm}$}
です。②で度の式から求めた値と同じです。

\noindent\textbf{度の式と同じになるわけ}\par
中心角 $a$° をラジアンにすると $\dfrac{\pi a}{180}$ です（④）。$s=r\theta$ に入れると、$r\times\dfrac{\pi a}{180}$。度の式は $2\pi r\times\dfrac{a}{360}=r\times\dfrac{2\pi a}{360}=r\times\dfrac{\pi a}{180}$ で、同じです。2つの式は、同じことを、ちがう単位で書いたものです。ラジアンのほうが、ずっと短く書けます。

\begin{column}{時計の針の先は、1分で何cm進む？}
\noindent\begin{minipage}[c]{0.66\linewidth}
長い針（分針）は、60分で1周（$2\pi$ rad）回ります。1分では $\dfrac{2\pi}{60}=\dfrac{\pi}{30}$ rad、約0.105\,rad です。長さ\qty{10}{cm}の長い針なら、その先は1分で、$s=r\theta=10\times\dfrac{\pi}{30}=\dfrac{\pi}{3}$ で、約\qty{1.05}{cm}進みます。

短い針（時針）は、12時間（720分）で1周なので、1分で $\dfrac{2\pi}{720}=\dfrac{\pi}{360}$ rad しか回りません。長さ\qty{6}{cm}の短い針の先は、1分で $6\times\dfrac{\pi}{360}=\dfrac{\pi}{60}$ で、約\qty{0.05}{cm}（\qty{0.5}{mm}）しか進まないので、動いているのが見えないのです。

ロンドンの時計台「ビッグ・ベン」の長い針は、長さが約\qty{4.2}{m}あります。その先は、1分で $4.2\times\dfrac{\pi}{30}\approx\qty{0.44}{m}$ も進みます。
\end{minipage}\hfill
\begin{minipage}[c]{0.3\linewidth}\centering
\begin{tikzpicture}[x=1mm, y=1mm]
  \bkclock{(0,0)}{15}{10}{10}
\end{tikzpicture}
\end{minipage}
\end{column}

\par\Needspace{16\baselineskip}\bkivcmidashi{おうぎ形の面積：細い三角形に分けて考える}
\noindent\begin{minipage}[c]{0.52\linewidth}
おうぎ形を、中心から細く細く切り分けます。1切れ1切れは、とても細いおうぎ形です。細いおうぎ形は、底辺が小さな弧、高さが半径 $r$ の、細長い三角形とほとんど同じです。

三角形の面積は、底辺×高さ÷2 です。どの三角形も高さは $r$ なので、ぜんぶ合わせると、
\[ \frac{1}{2}\times r\times(\text{底辺の合計}) \]
底辺の合計は、もとの弧の長さ $s$ です。
\end{minipage}\hfill
\begin{minipage}[c]{0.44\linewidth}\centering
\begin{tikzpicture}[x=1mm, y=1mm]
  % 中心角 60° のおうぎ形を、8つの細い三角形に分ける（半径 40 mm）
  \foreach \k in {0,...,7} {\pgfmathsetmacro\aa{60+\k*7.5}\pgfmathsetmacro\bb{60+(\k+1)*7.5}
    \pgfmathsetmacro\cc{mod(\k,2)==0 ? 30 : 12}
    \fill[bkOrange!\cc] (0,0) -- (\aa:40) -- (\bb:40) -- cycle;
    \draw[bkOrange!80!black, line width=0.4pt] (0,0) -- (\aa:40) -- (\bb:40) -- cycle;}
  \draw[bkBlue, line width=1.2pt] (60:40) arc (60:120:40);
  \fill[bkInk] (0,0) circle (0.8);
  \node[font=\sffamily\footnotesize, anchor=north] at (0,-1.4) {中心};
  \draw[bkInk, line width=0.5pt, {Stealth[length=2.5pt]}-{Stealth[length=2.5pt]}] ($(120:1)+(210:3)$) -- ($(120:39)+(210:3)$);
  \node[font=\sffamily\footnotesize, anchor=east] at ($(120:20)+(210:4.5)$) {高さ $r$};
  \node[font=\sffamily\footnotesize, text=bkBlue!80!black, anchor=south] at (0,41.2) {弧の長さ $s$};
\end{tikzpicture}
\end{minipage}

\bkivcline{\fbox{おうぎ形の面積　$S=\dfrac{1}{2}rs$}}
$S$ は面積によく使う文字です（英語の square、面積の単位の「平方」から）。この式は、三角形の面積「底辺×高さ÷2」と同じ形です。おうぎ形は、底辺が弧の長さ $s$、高さが半径 $r$ の三角形のようなもの、と考えればおぼえやすいでしょう。

\noindent\textbf{細かく分けるほど、正確になる}\par
三角形の底辺はまっすぐですが、弧は曲がっています。だから、細い三角形の面積の合計は、おうぎ形の面積より少しだけ小さくなります。でも、分ける数をふやすほど、そのちがいは小さくなります。半径\qty{6}{cm}、中心角60°のおうぎ形で、コンピュータに計算させると、次のようになります。

\begin{bkcenter}
\begin{bktable}{colspec={Q[l,m,wd=40mm] *{6}{Q[c,m,wd=12.5mm]}}, rowsep=1.5pt, colsep=2pt}
分ける数 & 1 & 2 & 4 & 8 & 16 & 100 \\
三角形の面積の合計（\unit{cm^2}） & 15.59 & 18.00 & 18.63 & 18.80 & 18.84 & 18.849 \\
\end{bktable}
\bktabcaption{$\dfrac{1}{2}rs=\dfrac{1}{2}\times6\times2\pi=6\pi\approx18.850$ に、どんどん近づく。}
\end{bkcenter}

このように「細かく分けて、たし合わせる」と、曲がった形の面積も求められます。この考え方は、4-5で円や球に、第5巻の積分（5-10）で、もっといろいろな形に使います。

\bkivcmidashi{角 $\theta$ で書いたおうぎ形の面積}
$S=\dfrac{1}{2}rs$ の $s$ に、$s=r\theta$ を入れます。
\bkivcline{\fbox{$S=\dfrac{1}{2}\times r\times r\theta=\dfrac{1}{2}r^2\theta$}\qquad（$\theta$ はラジアン）}
半径\qty{6}{cm}、中心角 $\dfrac{\pi}{3}$ なら、$S=\dfrac{1}{2}\times6^2\times\dfrac{\pi}{3}=\dfrac{36\pi}{6}=6\pi\approx\qty{18.85}{cm^2}$ です。

円全体は、中心角 $2\pi$ のおうぎ形です。$\theta=2\pi$ を入れると、
\bkivcline{$\dfrac{1}{2}r^2\times2\pi=\pi r^2$}
で、小学校で学んだ円の面積 $\pi r^2$ と、ちゃんと合っています。円の面積がなぜ $\pi r^2$ になるのかは、4-5で、円を細かく切って並べかえる方法でくわしく調べます。

\bksubsection{6}{円の接線は、半径に垂直}

\leavevmode\label{bk:4-2:sessen}円と直線の関係は、3通りあります。
\begin{bkfigure}
\begin{tikzpicture}[x=1mm, y=1mm]
  \foreach \xx/\yy/\lab in {0/-6/2点で交わる, 44/-14/1点でふれる, 88/-22/はなれている} {
    \begin{scope}[shift={(\xx,0)}]
      \draw[bkInk!70, line width=0.9pt] (0,0) circle (14);
      \fill[bkInk] (0,0) circle (0.7);
      \draw[bkOrange, line width=1.4pt] (-19,\yy) -- (19,\yy);
      \node[font=\sffamily\footnotesize, anchor=north] at (0,-25) {\lab};
    \end{scope}}
  \fill[bkInk] ({-sqrt(14*14-36)},-6) circle (1);
  \fill[bkInk] ({sqrt(14*14-36)},-6) circle (1);
  \fill[bkInk] (44,-14) circle (1);
\end{tikzpicture}
\bkfigcaption{直線（オレンジ）を下へずらしていくと、2点で交わる → 1点でふれる → はなれる、と変わる。}
\end{bkfigure}

まん中のように、円にちょうど1点だけでふれる直線を、円の\textbf{接線}（せっせん）といいます。\bkidx{せっせん}{接線}ふれている点を\textbf{接点}（せってん）といいます。接線は、円に「接する」直線です。丸いお皿のふちに定規をあてて、ぴったりふれさせたときの定規のへりが、接線です。

接線には、大切な性質があります。
\bkivcline{\textbf{接線は、接点で半径に垂直}}
なぜでしょう。2通りの方法で確かめます。

\bkivcstep{1}{方法1：2点で交わる直線を、近づけていく}
\noindent\begin{minipage}[c]{0.46\linewidth}\centering
\begin{tikzpicture}[x=1mm, y=1mm]
  % 中心O(0,0)、半径22。T は真下。P は $\angle\mathrm{TOP}$＝60° と 30° の所
  \draw[bkInk!60, line width=0.8pt] (0,0) circle (22);
  \fill[bkBlue!12] (0,0) -- (0,-22) -- (-30:22) -- cycle;
  \draw[bkBlue, line width=0.9pt] (0,0) -- (0,-22) -- (-30:22) -- cycle;
  \draw[bkGray, line width=0.8pt] (0,-22) -- ($(0,-22)!1.35!(-60:22)$);
  \draw[bkOrange, line width=1.5pt] (-26,-22) -- (30,-22);
  \fill[bkInk] (0,0) circle (0.9);
  \fill[bkInk] (0,-22) circle (0.9);
  \fill[bkInk] (-30:22) circle (0.9);
  \fill[bkInk] (-60:22) circle (0.9);
  \node[font=\sffamily\footnotesize, anchor=east] at (-1,0) {O};
  \node[font=\sffamily\footnotesize, anchor=north] at (-1.6,-23) {T};
  \node[font=\sffamily\footnotesize, anchor=west] at ($(-30:22)+(1,0.5)$) {P};
  \node[font=\sffamily\footnotesize, anchor=north west] at ($(-60:22)+(0.6,-0.3)$) {P$'$};
  \node[font=\sffamily\footnotesize, text=bkOrange!80!black, anchor=north] at (21,-23) {接線};
  \node[font=\sffamily\footnotesize, text=bkBlue!80!black] at (6.5,-9) {60°};
\end{tikzpicture}
\end{minipage}\hfill
\begin{minipage}[c]{0.5\linewidth}
円の中心をO、真下の点をTとします。円の上にもう1つの点Pをとり、直線TPを引きます。TPは、円と2点で交わる直線です。

三角形OTPで、OTもOPも半径なので、二等辺三角形です。だから、Tの角とPの角は等しくなります（4-1）。
\end{minipage}

\par\vspace{3pt}
Pを、円にそってTに近づけていきます（P→P$'$→……）。中心の角$\angle\mathrm{TOP}$は、だんだん小さくなります。三角形の内角の和は180°なので（4-1）、Tの角は $(180^{\circ}-\angle\mathrm{TOP})\div2$ です。

\begin{bkcenter}
\begin{bktable}{colspec={Q[l,m,wd=46mm] *{5}{Q[c,m,wd=13mm]}}, rowsep=1.5pt}
中心の角$\angle\mathrm{TOP}$ & 60° & 30° & 10° & 2° & 0.2° \\
Tの角（直線TPと半径OTの角） & 60° & 75° & 85° & 89° & 89.9° \\
\end{bktable}
\bktabcaption{Pを近づけるほど、直線TPと半径OTの角は90°に近づく。}
\end{bkcenter}

Pが近づくほど、直線TPは、Tでふれる1本の直線（接線）に近づいていきます。そのとき、半径OTとの角は90°に近づきます。だから、接線と半径の角は90°、つまり垂直です。5-3では、この「近づけていく」考え方で、曲線の接線を求めます。

\bkivcstep{2}{方法2：タイヤと道路を、折り目で考える}
\noindent\begin{minipage}[c]{0.56\linewidth}
平らな道路にのっているタイヤを、横から見ます。タイヤ（円）は、道路（直線）に1点だけでふれています。道路は、タイヤの接線です。

タイヤの中心（車軸）から、道路にまっすぐ垂直に下ろした線を引きます。この線で、紙を折ることを考えます。円は、中心を通る線で折ると、左右がぴったり重なります。道路も、この折り目に垂直なので、折ると自分自身に重なります。
\end{minipage}\hfill
\begin{minipage}[c]{0.4\linewidth}\centering
\begin{tikzpicture}[x=1mm, y=1mm]
  \draw[bkGray, line width=1.2pt] (-24,0) -- (24,0);
  \fill[bkGray!25] (-24,-3) rectangle (24,0);
  \draw[bkInk!80, line width=3pt] (0,15.5) circle (15);
  \draw[bkInk!40, line width=0.8pt] (0,15.5) circle (9);
  \fill[bkInk!60] (0,15.5) circle (2);
  \draw[bkBlue, line width=0.9pt, dash pattern=on 2.5pt off 1.5pt] (0,35) -- (0,-5);
  \fill[bkOrange] (0,0) circle (1.2);
  \bkivcright{(0,0)}{0}{90}
  \node[font=\sffamily\footnotesize, text=bkBlue!80!black, anchor=west] at (0.8,34) {折り目};
  \node[font=\sffamily\footnotesize, anchor=north east] at (25,-3.5) {道路（接線）};
  \node[font=\sffamily\footnotesize, text=bkOrange!80!black, anchor=north east] at (-1.4,-3.5) {接点};
\end{tikzpicture}
\end{minipage}

\par\vspace{3pt}
もし、タイヤと道路のふれる点が、折り目より右にずれていたとします。折ると、その点は折り目の左側にうつります。円も道路も自分自身に重なるので、うつった点も、円と道路がふれる点です。すると、ふれる点が左右に2つできてしまいます。でも、接線がふれる点は1つだけです。だから、ふれる点（接点）は、折り目の上にあるはずです。つまり接点は、中心の真下です。中心と接点を結んだ半径は折り目そのもので、道路に垂直です。

\begin{taisetsu}
円の接線は、円に1点だけでふれる直線。\\
接線は、ふれた点（接点）で半径に垂直。
\end{taisetsu}

この性質を使うと、円の上の点Tでの接線は、「半径OTに垂直で、Tを通る直線」として引けます。

\bksubsection{7}{垂直二等分線と、円の中心}

\leavevmode\label{bk:4-2:gaishin}紙に2点A、Bをかき、AとBがぴったり重なるように紙を折ってみましょう。折り目は、線分ABのまん中を通り、ABに垂直です。この直線を、線分ABの\textbf{垂直二等分線}といいます（「垂直」で、「2つの等しい長さに分ける」線）。

\noindent\begin{minipage}[c]{0.44\linewidth}\centering
\begin{tikzpicture}[x=1mm, y=1mm]
  \draw[bkInk, line width=1pt] (-18,0) -- (18,0);
  \fill[bkInk] (-18,0) circle (1); \fill[bkInk] (18,0) circle (1);
  \node[font=\sffamily\footnotesize, anchor=north] at (-18,-1.6) {A};
  \node[font=\sffamily\footnotesize, anchor=north] at (18,-1.6) {B};
  \draw[bkBlue, line width=1.1pt] (0,-8) -- (0,32);
  \bkivcright{(0,0)}{0}{90}
  \foreach \yy/\c in {12/bkOrange, 26/bkPurple} {
    \draw[\c, line width=0.8pt, dash pattern=on 2pt off 1.2pt] (-18,0) -- (0,\yy) -- (18,0);
    \fill[\c] (0,\yy) circle (1);}
  \node[font=\sffamily\footnotesize, anchor=west] at (1.2,13.5) {P};
  \node[font=\sffamily\footnotesize, anchor=west] at (1.2,27.5) {Q};
  \node[font=\sffamily\footnotesize, text=bkBlue!80!black, anchor=north west] at (1.4,-2) {垂直二等分線};
\end{tikzpicture}
\end{minipage}\hfill
\begin{minipage}[c]{0.52\linewidth}
垂直二等分線の上に点Pをとります。紙を折るとAとBが重なるので、PからAまでの長さと、PからBまでの長さは同じです。点Qでも同じです。

\par\addvspace{3pt}
\textbf{垂直二等分線の上の点は、AからもBからも同じ長さ。}\par\noindent
反対に、AからもBからも同じ長さの点は、垂直二等分線の上にあります。
\end{minipage}

\bkivcmidashi{弦の垂直二等分線は、円の中心を通る}
円の弦ABを考えます。円の中心Oから、AまでもBまでも、長さは半径で同じです。だから、Oは弦ABの垂直二等分線の上にあります。

\noindent\begin{minipage}[c]{0.54\linewidth}
どの弦の垂直二等分線も、中心を通ります。だから、弦を2本引いて、それぞれの垂直二等分線を引けば、2本の交わる点が円の中心です。

たとえば、割れてしまった丸いお皿のかけらから、もとのお皿の中心と大きさがわかります。新しいお皿を注文したり、こわれた部品を作り直したりするときに使う方法です。
\end{minipage}\hfill
\begin{minipage}[c]{0.42\linewidth}\centering
\begin{tikzpicture}[x=1mm, y=1mm]
  % お皿のかけら（中心 (0,0)、半径 24 の円の一部）
  \fill[bkYellow!35] (200:24) arc (200:330:24) -- (330:17) -- (300:12) -- (250:15) -- (200:17) -- cycle;
  \draw[bkYellow!50!black, line width=1.2pt] (200:24) arc (200:330:24);
  \draw[bkInk!40, line width=0.5pt, dash pattern=on 1.5pt off 1.5pt] (-24,0) arc (180:200:24);
  \draw[bkInk!40, line width=0.5pt, dash pattern=on 1.5pt off 1.5pt] (330:24) arc (330:360:24);
  \draw[bkOrange, line width=1pt] (205:24) -- (260:24);
  \draw[bkPurple, line width=1pt] (275:24) -- (325:24);
  \draw[bkOrange!80!black, line width=0.7pt, dash pattern=on 2pt off 1pt] (232.5:27) -- (0,0);
  \draw[bkPurple!80!black, line width=0.7pt, dash pattern=on 2pt off 1pt] (300:27) -- (0,0);
  \fill[bkInk] (0,0) circle (1.1);
  \node[font=\sffamily\footnotesize, anchor=south] at (0,1.4) {中心};
\end{tikzpicture}
\end{minipage}

\bkivcmidashi{3つの点を通る円：外接円と外心}
三角形ABCの3つの頂点を、ぜんぶ通る円をかきたいとします。その円の中心は、AからもBからもCからも同じ長さの点です。AとBから同じ長さなので、辺ABの垂直二等分線の上にあります。AとCから同じ長さなので、辺ACの垂直二等分線の上にもあります。だから、2本の垂直二等分線の交わる点が中心です。

三角形の3つの頂点を通る円を\textbf{外接円}（がいせつえん）、その中心を\textbf{外心}（がいしん）といいます。\bkidx{がいしん}{外心}4-7では、外接円の半径を使って、三角形の辺と角の関係を調べます。

\bksubsection{8}{円周角は、中心角の半分}

\leavevmode\label{bk:4-2:enshuu}円の上に2点A、Bと、もう1つの点Pをとります。角$\angle\mathrm{APB}$を、弧ABに対する\textbf{円周角}（えんしゅうかく）といいます。\bkidx{えんしゅうかく}{円周角}頂点Pが円の上にある角です。

\begin{bkfigure}
\begin{tikzpicture}[x=1mm, y=1mm]
  % 左：中心角 80°、円周角は P の場所によらず 40°（A は 230°、B は 310°）
  \begin{scope}[shift={(0,0)}]
    \draw[bkInk!60, line width=0.8pt] (0,0) circle (22);
    \draw[bkOrange, line width=1.8pt] (230:22) arc (230:310:22);
    \fill[bkTheme!45] (0,0) -- (230:22) -- (310:22) -- cycle;
    \draw[bkInk, line width=0.7pt] (230:22) -- (0,0) -- (310:22);
    \foreach \p/\c in {40/bkBlue, 100/bkPurple, 155/bkRed} {
      \draw[\c, line width=0.8pt] (230:22) -- (\p:22) -- (310:22);
      \fill[\c] (\p:22) circle (1);}
    \draw[bkInk, line width=0.5pt] ($(0,0)+(230:4.5)$) arc (230:310:4.5);
    \node[font=\sffamily\footnotesize] at (270:12.5) {80°};
    \fill[bkInk] (0,0) circle (0.8);
    \node[font=\sffamily\footnotesize, anchor=north east] at (230:22.6) {A};
    \node[font=\sffamily\footnotesize, anchor=north west] at (310:22.6) {B};
    \node[font=\sffamily\footnotesize, text=bkBlue!80!black, anchor=south west] at (40:22.4) {40°};
    \node[font=\sffamily\footnotesize, text=bkPurple!80!black, anchor=south] at (100:22.6) {40°};
    \node[font=\sffamily\footnotesize, text=bkRed!80!black, anchor=south east] at (155:22.4) {40°};
  \end{scope}
  % 右：直径の円周角は 90°
  \begin{scope}[shift={(66,0)}]
    \draw[bkInk!60, line width=0.8pt] (0,0) circle (22);
    \draw[bkOrange, line width=1.2pt] (-22,0) -- (22,0);
    \foreach \p/\c in {55/bkBlue, 125/bkPurple} {
      \draw[\c, line width=0.8pt] (-22,0) -- (\p:22) -- (22,0);
      \fill[\c] (\p:22) circle (1);}
    \bkivcright{(55:22)}{207.5}{297.5}
    \fill[bkInk] (0,0) circle (0.8);
    \node[font=\sffamily\footnotesize, anchor=north] at (0,-1.2) {O};
    \node[font=\sffamily\footnotesize, anchor=east] at (-23,0) {A};
    \node[font=\sffamily\footnotesize, anchor=west] at (23,0) {B};
    \node[font=\sffamily\footnotesize, text=bkBlue!80!black, anchor=south west] at (55:22.4) {90°};
    \node[font=\sffamily\footnotesize, text=bkPurple!80!black, anchor=south east] at (125:22.4) {90°};
    \node[font=\sffamily\footnotesize, text=bkOrange!80!black, anchor=north] at (11,-1.2) {直径};
  \end{scope}
\end{tikzpicture}
\bkfigcaption{左：中心角$\angle\mathrm{AOB}$が80°のとき、円周角は、Pがどこにあっても40°。右：ABが直径なら、円周角は90°。}
\end{bkfigure}

分度器ではかると、Pを円の上のどこに動かしても（弧ABの上に入らなければ）、$\angle\mathrm{APB}$は同じ40°です。中心角$\angle\mathrm{AOB}$の80°の、ちょうど半分です。
\bkivcline{\fbox{円周角＝中心角の半分}}

\par\Needspace{20\baselineskip}\bkivcmidashi{なぜ半分になるの？}
\noindent\begin{minipage}[c]{0.6\linewidth}
Pから中心Oを通る直径PCを引いた場合で考えます。三角形OPBは、OPもOBも半径なので、二等辺三角形です。2つの底角は等しいので、どちらも ● とします（4-1）。

$\angle\mathrm{BOC}$は、三角形OPBの外角です。外角は、となりにない2つの内角の和なので（4-1）、$\angle\mathrm{BOC}$＝●＋●＝2×●。中心角$\angle\mathrm{BOC}$は、円周角$\angle\mathrm{BPC}$（＝●）の2倍です。ほかの場所にPがあるときも、このような形を2つ組み合わせて、同じことが確かめられます。
\end{minipage}\hfill
\begin{minipage}[c]{0.36\linewidth}\centering
\begin{tikzpicture}[x=0.82mm, y=0.82mm]
  % P は 110°、C は 290°（PC が直径）、B は 350°
  \draw[bkInk!60, line width=0.8pt] (0,0) circle (21);
  \fill[bkBlue!10] (110:21) -- (0,0) -- (350:21) -- cycle;
  \draw[bkInk, line width=0.8pt] (110:21) -- (290:21);
  \draw[bkBlue, line width=0.9pt] (110:21) -- (350:21) -- (0,0);
  \fill[bkInk] (0,0) circle (0.8);
  \foreach \a in {110,290,350} \fill[bkInk] (\a:21) circle (0.9);
  \node[font=\sffamily\footnotesize, anchor=south east] at (110:21.4) {P};
  \node[font=\sffamily\footnotesize, anchor=north] at (290:21.8) {C};
  \node[font=\sffamily\footnotesize, anchor=west] at (350:21.6) {B};
  \node[font=\sffamily\footnotesize, anchor=east] at (-1.2,1.2) {O};
  \node[font=\sffamily\footnotesize, text=bkBlue!80!black] at ($(110:21)+(-45:5.5)$) {●};
  \node[font=\sffamily\footnotesize, text=bkBlue!80!black] at ($(350:21)+(155:5.5)$) {●};
  \draw[bkOrange, line width=0.8pt] (290:5) arc (290:350:5);
  \node[font=\sffamily\footnotesize, text=bkOrange!80!black] at (320:8.6) {2●};
\end{tikzpicture}
\end{minipage}

\bkivcmidashi{直径に対する円周角は90°}
ABが直径のとき、中心角$\angle\mathrm{AOB}$は、まっすぐな角180°です。円周角はその半分なので、
\bkivcline{\textbf{直径に対する円周角は90°}}
です。直径の両はしと、円の上の点を結ぶと、いつも直角三角形ができます。大工さんは、この性質を使って、直角三角定規（さしがね）の直角の頂点を円にあて、2本の辺が円と交わる2点を結んで、直径（と中心）を見つけます。

\bksubsection{9}{電卓と表計算の DEG・RAD}

\leavevmode\label{bk:4-2:degrad}角を使う計算は、電卓や表計算ですることが多くなります。そのとき、角を度で入れるのか、ラジアンで入れるのかに注意が必要です。

\noindent\begin{minipage}[c]{0.62\linewidth}
\textbf{関数電卓}\par
角を使う電卓（関数電卓）には、角の単位を切りかえるボタンがあります。画面に DEG と出ていれば度、RAD と出ていればラジアンで計算します。DEG は degree（度）、RAD は radian（ラジアン）の頭の3文字です。

4-6からは、角を入れると数が出てくる $\sin$（サイン）などの計算が出てきます。そのとき、DEG と RAD をまちがえると、まったくちがう答えが出ます。計算をはじめる前に、画面の DEG・RAD を必ず確かめましょう。
\end{minipage}\hfill
\begin{minipage}[c]{0.34\linewidth}\centering
\begin{tikzpicture}[x=1mm, y=1mm]
  % 第14回 B2：DRG キーを広くして字をキーの中におさめ、赤い丸の中心を DRG キーの中心 (29.8, 27.5) にそろえる
  \draw[bkInk!70, line width=1pt, rounded corners=3pt, fill=bkPaper] (0,0) rectangle (37.6,48);
  \draw[bkInk!50, line width=0.6pt, fill=white] (3,34) rectangle (34.6,45);
  \node[font=\sffamily\footnotesize\bfseries, text=bkRed!80!black, anchor=north west, inner sep=1pt] at (4,44.4) {RAD};
  \node[font=\ttfamily\small, anchor=south east] at (33.6,35) {1.047197551};
  \foreach \x/\t in {3/sin, 10.3/cos, 17.6/tan} {\draw[bkInk!40, fill=bkTheme!12, rounded corners=1pt] (\x,25) rectangle ++(6.6,5);
    \node[font=\sffamily\footnotesize] at (\x+3.3,27.5) {\t};}
  \draw[bkInk!40, fill=bkTheme!12, rounded corners=1pt] (25.3,25) rectangle (34.3,30);
  \node[font=\sffamily\footnotesize] at (29.8,27.5) {DRG};
  \foreach \i in {0,1,2,3} \foreach \j in {0,1,2} {\draw[bkInk!40, fill=white, rounded corners=1pt] (3+\i*8.1,3+\j*6.8) rectangle ++(7.3,5);}
  \draw[bkRed!80!black, line width=0.8pt] (29.8,27.5) circle (5.3);
\end{tikzpicture}
\end{minipage}

\bkivcmidashi{表計算は、いつもラジアン}
表計算ソフトの \texttt{SIN} などの関数は、いつもラジアンで計算します。DEG のような切りかえはありません。度で計算したいときは、関数 \texttt{RADIANS} で、先にラジアンに直します。

\begin{bkcenter}
\bksheet{3}{30}{%
1 & 度 & \texttt{=RADIANS(A2)} & \texttt{=DEGREES(B2)} \\
2 & 30 & 0.5236 & 30 \\
3 & 60 & 1.0472 & 60 \\
4 & 90 & 1.5708 & 90 \\
5 & 180 & 3.1416 & 180 \\
6 & 360 & 6.2832 & 360 \\}
\bktabcaption{\texttt{RADIANS} は度をラジアンに、\texttt{DEGREES} はラジアンを度に直す。180°は3.1416（$\pi$）。}
\end{bkcenter}

表計算のセルに \texttt{=PI()} と入れると、$\pi=3.14159\cdots$ が出ます。\texttt{=RADIANS(A2)} という式は、\texttt{=A2*PI()/180} と同じ計算をしています。④で作った「$\times\dfrac{\pi}{180}$」です。

たとえば、4-6で学ぶ「30°のサイン」は0.5です。表計算で \texttt{=SIN(30)} と入れると、30ラジアンのサインを計算してしまい、$-0.988\cdots$ というちがう数が出ます。30°なら、\texttt{=SIN(RADIANS(30))} と入れると、0.5 が出ます。プログラム（2-19）でも、ふつうはラジアンで計算します。Python なら \texttt{math.radians(30)} で度をラジアンに直せます。

%% ---------------- ⑩ 極座標（第14回 U-P：三角比の前に、距離と向きで場所を言う） ----------------
\bksubsection{10}{距離と向きで場所を言う：極座標}

\leavevmode\label{bk:4-2:kyoku}夜の海で、灯台のレーダーが船を見つけました。レーダーは、アンテナをぐるぐる回しながら電波を出します。電波が船に当たってはね返り、もどってくるまでの時間から、船までの距離がわかります。そのとき、アンテナがどちらを向いていたかで、船のいる向きもわかります。だから、レーダーの画面では、船の場所が「中心から何km、どの向き」という2つの数で表されます。

\begin{bkfigure}
\begin{tikzpicture}[x=1mm, y=1mm]
  % 左：レーダーの画面（1 km を 6 mm。円は 1 km ごと、線は 30°ごと）
  \bkivcpolar{6}{4}
  \draw[bkOrange, line width=1.4pt] (0,0) -- (60:18);
  \draw[bkOrange!80!black, line width=0.6pt, -{Stealth[length=2.5pt]}] (5,0) arc (0:60:5);
  \node[font=\sffamily\footnotesize, fill=white, inner sep=0.8pt] at (25:9.5) {60°};
  \fill[bkOrange!80!black] (60:18) circle (1.3);
  \node[font=\sffamily\footnotesize, fill=white, inner sep=0.8pt, anchor=west] at ($(60:18)+(2,-0.6)$) {船A};
  \node[font=\sffamily\footnotesize, fill=white, inner sep=0.8pt, rotate=60] at ($(60:10.5)+(150:3.3)$) {\qty{3}{km}};
  % 右：ロボットのうで（関節が1つ。肩 S、長さ 30 cm を 30 mm、水平から 40°）
  \begin{scope}[shift={(60,-12)}]
    \fill[bkGray!25] (-10,-9.5) rectangle (34,-8);
    \draw[bkGray, line width=0.8pt] (-10,-8) -- (34,-8);
    \fill[bkGray!60] (-6,-8) [rounded corners=0.8mm] -- (-6,-3) -- (6,-3) [sharp corners] -- (6,-8) -- cycle;
    \fill[bkGray!80] (-3.5,-3) -- (3.5,-3) -- (2,-1) -- (-2,-1) -- cycle;
    \draw[bkGray, line width=0.6pt, dash pattern=on 2pt off 1.5pt] (0,0) -- (30,0);
    \node[font=\sffamily\footnotesize, anchor=west] at (31,0) {0°};
    \draw[bkBlue, line width=3.2pt, line cap=round] (0,0) -- (40:30);
    \begin{scope}[shift={(40:30)}, rotate=40]
      \draw[bkInk, line width=1pt, line cap=round] (0.4,-1.7) -- (0.4,1.7) (0.4,-1.7) -- (2.4,-1.7) (0.4,1.7) -- (2.4,1.7);
    \end{scope}
    \filldraw[fill=white, draw=bkBlue!70!black, line width=0.8pt] (0,0) circle (1.5);
    \draw[bkInk, line width=0.5pt] (9,0) arc (0:40:9);
    \node[font=\sffamily\footnotesize] at (20:14) {40°};
    \node[font=\sffamily\footnotesize, text=bkBlue!80!black, rotate=40] at ($(40:15)+(130:4)$) {\qty{30}{cm}};
    \node[font=\sffamily\footnotesize, anchor=south, align=center] at ($(40:30)+(-1,4.5)$) {手};
  \end{scope}
\end{tikzpicture}
\bkfigcaption{左：レーダーの画面。円は中心から1・2・3・4\,km、線は30°ごと。船Aは、中心から\qty{3}{km}、0°の線から左回りに60°の向き。右：ロボットのうで（4-0）。手は、肩から\qty{30}{cm}、水平な線から40°の向き。}
\end{bkfigure}

ロボットのうでも同じです（4-0）。関節が1つで、長さ\qty{30}{cm}のうでを、水平な線から40°だけ回して上げると、手の場所は「肩から\qty{30}{cm}、向き40°」で1つに決まります。ロボットのモーターが決めるのは回す角度で、うでの長さは変わりません。だから、ロボットにとっては「距離と向き」で場所を言うのが、いちばん自然なのです。

\bkivcmidashi{極座標の約束}
このように、中心からの\textbf{距離}と、中心から見た\textbf{向き}の2つの数で場所を言う方法を、次のように決めます。
\begin{enumerate}[label={\bkcircled[bkOrange]{\arabic*}}, leftmargin=2.4\zw, itemsep=1pt, topsep=3pt]
\item 中心にする点Oを決める。3-2の座標の原点にあたる点です。
\item Oから右へのびる向き（$x$ 軸の正の向き）を、角をはかりはじめる向き（0°）にする。
\item 点Pについて、OからPまでの距離を $r$、0°の向きから\textbf{左回り}（時計の針と反対の向き）にはかったPの向きの角を $\theta$ とする。
\end{enumerate}
2つの数の組 $(r,\ \theta)$ を、点Pの\textbf{極座標}（きょくざひょう）といいます。\bkidx{きょくざひょう}{極座標}「極」は、中心の点Oのことです。レーダーの船Aの極座標は $(3,\ 60°)$（距離の単位は km）、ロボットの手の極座標は $(30,\ 40°)$（距離の単位は cm）です。

\bkivcmidashi{なぜ左回りにはかるの？}
\noindent\begin{minipage}[c]{0.6\linewidth}
3-2の座標では、右が $x$ 軸の正の向き、上が $y$ 軸の正の向きでした。右（0°）から上へ、直角だけ回る向きは、左回りです。だから、右から左回りにはかると、真上がちょうど90°、真左が180°、真下が270°になって、座標の軸と角がきれいにそろいます。数学では、この左回りを「正の向き」と決めています（右回りの角は 4-8 で考えます）。
\end{minipage}\hfill
\begin{minipage}[c]{0.36\linewidth}\centering
\begin{tikzpicture}[x=1mm, y=1mm]
  \draw[bkGray, line width=0.6pt, -{Stealth[length=4pt]}] (-17,0) -- (17,0) node[anchor=west, font=\sffamily\footnotesize, text=bkInk] {$x$};
  \draw[bkGray, line width=0.6pt, -{Stealth[length=4pt]}] (0,-17) -- (0,17) node[anchor=south, font=\sffamily\footnotesize, text=bkInk] {$y$};
  \draw[bkOrange, line width=1.1pt, -{Stealth[length=4pt]}] (8,0) arc (0:340:8);
  \node[font=\sffamily\footnotesize, fill=white, inner sep=0.6pt] at (12.5,-2.6) {0°};
  \node[font=\sffamily\footnotesize, fill=white, inner sep=0.6pt] at (3.6,12) {90°};
  \node[font=\sffamily\footnotesize, fill=white, inner sep=0.6pt] at (-12.8,-2.6) {180°};
  \node[font=\sffamily\footnotesize, fill=white, inner sep=0.6pt] at (4.6,-12.5) {270°};
  \fill[bkInk] (0,0) circle (0.8);
  \node[font=\sffamily\footnotesize, anchor=north east, inner sep=1pt] at (-0.6,-0.6) {O};
\end{tikzpicture}
\end{minipage}

\par\vspace{4pt}\bkivcmidashi{$(x,\ y)$ と $(r,\ \theta)$：同じ点の、2つの名前}
3-2の座標 $(x,\ y)$ は、「右へいくつ、上へいくつ」で場所を言いました。極座標 $(r,\ \theta)$ は、「どの向きへ、どれだけ」で場所を言います。同じ点を、2通りの名前でよべるのです。軸の上の点なら、すぐに行き来できます。
\par\vspace{3pt}\noindent\begin{minipage}[c]{0.5\linewidth}\centering
\begin{bktable}{colspec={Q[l,m,wd=24mm] Q[c,m,wd=17mm] Q[c,m,wd=19mm]}, rowsep=1.2pt, colsep=3pt, row{1}={font=\sffamily\small}}
点の場所 & $(x,\ y)$ & $(r,\ \theta)$ \\
Oから右へ3 & $(3,\ 0)$ & $(3,\ 0°)$ \\
Oから上へ2 & $(0,\ 2)$ & $(2,\ 90°)$ \\
Oから左へ4 & $(-4,\ 0)$ & $(4,\ 180°)$ \\
Oから下へ1 & $(0,\ -1)$ & $(1,\ 270°)$ \\
\end{bktable}
\bktabcaption{軸の上の点の、2つの名前}
\end{minipage}\hfill
\begin{minipage}[c]{0.46\linewidth}
たとえば $(-4,\ 0)$ は、Oから左へ4の点です。左は180°の向きなので、極座標は $(4,\ 180°)$ です。$r$ は距離なので、マイナスにはなりません。向きのちがいは、$\theta$ が受けもちます。
\end{minipage}

\par\vspace{4pt}
ななめの向きの点、たとえば船A $(3,\ 60°)$ は、中心から右へいくつ、上へいくつの所にあるでしょう。方眼紙に図を大きくかいて読みとると、右へ約1.5、上へ約2.6です。これを、図をかかずに計算で求める道具が、4-6で作る $\sin$（サイン）・$\cos$（コサイン）です。4-6では、極座標 $(r,\ \theta)$ を $(x,\ y)$ に直す式を作ります。いまは、「同じ点を、$(x,\ y)$ でも $(r,\ \theta)$ でも言える」ことがわかれば十分です。

\bkivcmidashi{角は、度でもラジアンでもよい}
$\theta$ は角なので、度でもラジアンでも書けます。60°＝$\dfrac{\pi}{3}$ rad なので（④）、船Aは $(3,\ 60°)$ とも、$\left(3,\ \dfrac{\pi}{3}\right)$ とも書けます。ラジアンには単位をつけないことが多いので（④のせいかくに言うと）、「°」がついていない角はラジアンです。ふつうは、$\theta$ を0°以上360°未満（ラジアンなら0以上 $2\pi$ 未満）で答えます。360°より大きい角や、マイナスの角は、4-8で考えます。

\begin{seikaku}
\textbf{$r$ と $\theta$ のきまり}\par
$r$ は radius（半径）の頭文字です。中心からの距離なので、0以上の数です。3-15の相関係数の $r$ とは、別の意味です（同じ文字が、場面によってちがう意味で使われます）。中心の点Oは $r=0$ で、向きは決まりません（どの向きに0だけ進んでも、Oのままだから）。極座標は、$(r,\ \theta)$ の順に書きます。$(60°,\ 3)$ のように逆にしないこと。
\end{seikaku}


\begin{taisetsu}
円周＝$2\pi r$。弧の長さとおうぎ形の面積は、中心角に比例する。\par\vspace{2pt}
弧度法：角 $\theta$（rad）＝弧の長さ÷半径。半径1の円なら、角と弧の長さは同じ数。\\
1周は $2\pi$\,rad（360°）。$\pi$\,rad＝180°\\
$1^{\circ}=\dfrac{\pi}{180}$\,rad\qquad $\qty{1}{rad}=\dfrac{180}{\pi}^{\circ}\approx57.3^{\circ}$\par\vspace{2pt}
弧の長さ＝半径×角\quad $s=r\theta$\qquad おうぎ形の面積\quad $S=\dfrac{1}{2}rs=\dfrac{1}{2}r^2\theta$\par\vspace{2pt}
円周角は中心角の半分。直径に対する円周角は90°。\par\vspace{2pt}
電卓は DEG（度）と RAD（ラジアン）を確かめる。表計算の \texttt{SIN} などはラジアン。\par\vspace{2pt}
極座標 $(r,\ \theta)$：中心からの距離 $r$ と、$x$ 軸の正の向きから左回りにはかった向きの角 $\theta$ で、場所を言う。
\end{taisetsu}

\begin{kotoba}
\small
\bkword{弦・弧}{円の上の2点を結ぶ線分が弦。円の上の2点の間の曲がった部分が弧。}
\bkword{おうぎ形・中心角}{2本の半径と弧でかこまれた形がおうぎ形。2本の半径の間の角が中心角。}
\bkword{弧度法・ラジアン}{角を「弧の長さ÷半径」で表す方法が弧度法。その単位がラジアン（rad）。}
\bkword{接線・接点}{円に1点だけでふれる直線が接線。ふれる点が接点。接線は接点で半径に垂直。}
\bkword{垂直二等分線}{線分のまん中を通り、線分に垂直な直線。両はしから同じ長さの点の集まり。}
\bkword{外接円・外心}{三角形の3つの頂点を通る円が外接円。その中心が外心。}
\bkword{円周角}{円の上に頂点がある角。同じ弧に対する円周角は、中心角の半分。}
\bkword{極座標}{中心からの距離 $r$ と、$x$ 軸の正の向きから左回りにはかった角 $\theta$ の組 $(r,\ \theta)$ で、点の場所を表す方法。}
\end{kotoba}

\bkkigocards{\bkkigocard{$\theta$}{シータ}{角の大きさ。度でもラジアンでも使う（微分・積分ではラジアン）}{theta}{4-2}}
{\bkkigocard{rad}{ラジアン}{弧度法の角の単位。弧の長さ÷半径。$\pi$\,rad＝180°}{radian}{4-2}}

%% ---------------- 例題 ----------------
\bksection{reidai}{例題でたしかめよう}

\begin{reidai}[度で、弧の長さとおうぎ形の面積を求める]
半径\qty{6}{cm}、中心角60°のおうぎ形の、（1）\nobreak 弧の長さと、（2）\nobreak 面積を求めましょう。答えは $\pi$ を使って表し、小数（約）でも答えましょう。\par
\kangaekata （1）中心角60°は、円全体（360°）の $\dfrac{60}{360}$。弧の長さも円周の $\dfrac{60}{360}$ 倍です。（2）弧の長さ $s$ がわかれば、$S=\dfrac{1}{2}rs$ が使えます。\par
\shiki （1）円周は $2\pi\times6=12\pi$\par
\hspace*{4.3\zw}弧の長さ＝$12\pi\times\dfrac{60}{360}=12\pi\times\dfrac{1}{6}=2\pi$\par
\hspace*{4.3\zw}$2\pi\approx2\times3.14=6.28$\par
\hspace*{2.6\zw}（2）$S=\dfrac{1}{2}\times6\times2\pi=6\pi$\par
\hspace*{4.3\zw}$6\pi\approx6\times3.14=18.84$。くわしい $\pi$ で計算すると18.85\par
\kotae （1）\nobreak $2\pi$\,cm（約\qty{6.28}{cm}）\quad（2）\nobreak $6\pi$\,\unit{cm^2}（約\qty{18.8}{cm^2}）
\end{reidai}

\begin{jibun}
半径\qty{4}{cm}、中心角135°のおうぎ形の、（1）\nobreak 弧の長さと、（2）\nobreak 面積を求めましょう。$\pi$ を使って表し、小数（約）でも答えましょう。
\bkjans{（1）\nobreak 円周 $8\pi$ の $\dfrac{135}{360}=\dfrac{3}{8}$ で $3\pi$\,cm（約\qty{9.42}{cm}）。（2）\nobreak $\dfrac{1}{2}\times4\times3\pi=6\pi$\,\unit{cm^2}（約\qty{18.8}{cm^2}）。}
\end{jibun}
\bkivcpoint{度で求めるときは、① 円周 $2\pi r$ を出す、② 「中心角÷360」をかける、の2段で考えます。面積は、弧の長さがわかれば $S=\dfrac{1}{2}rs$ を使うと早く求められます。}

\begin{reidai}[度とラジアンを換算する]
（1）\nobreak 45°と150°を、ラジアンで表しましょう。\par
（2）\nobreak $\dfrac{2\pi}{3}$ rad と 2\,rad を、度で表しましょう（2\,rad は小数第1位まで）。\par
\kangaekata 「単位つきの1」$\dfrac{\pi\ \text{rad}}{180^{\circ}}$ または $\dfrac{180^{\circ}}{\pi\ \text{rad}}$ を、消したい単位が分母にくるようにかけます。\par
\begin{bkcenter}
\begin{tikzpicture}[x=1mm, y=1mm]
  % 度とラジアンの2つの目もり（180°を 108 mm）
  \draw[bkInk, line width=0.8pt] (0,0) -- (108,0);
  \foreach \d/\r in {0/0, 30/\frac{\pi}{6}, 45/\frac{\pi}{4}, 60/\frac{\pi}{3}, 90/\frac{\pi}{2}, 120/\frac{2\pi}{3}, 150/\frac{5\pi}{6}, 180/\pi} {
    \draw[bkInk, line width=0.6pt] ({\d*0.6},-1.4) -- ({\d*0.6},1.4);
    \node[font=\sffamily\footnotesize, anchor=south] at ({\d*0.6},2) {\d°};
    \node[font=\footnotesize, anchor=north] at ({\d*0.6},-2) {$\displaystyle\r$};}
  \draw[bkOrange, line width=1.2pt] (68.75,-1.8) -- (68.75,1.8);
  \node[font=\sffamily\footnotesize, text=bkOrange!80!black, anchor=north] at (68.75,-11.5) {2 rad（約114.6°）};
  \node[font=\sffamily\footnotesize, anchor=east] at (-3,3.5) {度};
  \node[font=\sffamily\footnotesize, anchor=east] at (-3,-5) {rad};
\end{tikzpicture}
\end{bkcenter}

\shiki （1）$45^{\circ}\times\dfrac{\pi\ \text{rad}}{180^{\circ}}=\dfrac{45\pi}{180}\ \text{rad}=\dfrac{\pi}{4}\ \text{rad}$\quad（45でわって約分）\par
\par\vspace{2pt}\hspace*{4.3\zw}$150^{\circ}\times\dfrac{\pi\ \text{rad}}{180^{\circ}}=\dfrac{150\pi}{180}\ \text{rad}=\dfrac{5\pi}{6}\ \text{rad}$\quad（30でわって約分）\par
\par\vspace{2pt}\hspace*{2.6\zw}（2）$\dfrac{2\pi}{3}\ \text{rad}\times\dfrac{180^{\circ}}{\pi\ \text{rad}}=\dfrac{2\times180^{\circ}}{3}=120^{\circ}$\quad（$\pi$ が約分される）\par
\par\vspace{2pt}\hspace*{4.3\zw}$2\ \text{rad}\times\dfrac{180^{\circ}}{\pi\ \text{rad}}=\dfrac{360^{\circ}}{\pi}=360^{\circ}\div3.14159\cdots\approx114.6^{\circ}$\par
\kotae （1）\nobreak $\dfrac{\pi}{4}$ rad、$\dfrac{5\pi}{6}$ rad\quad（2）\nobreak 120°、約114.6°
\end{reidai}

\begin{jibun}
（1）\nobreak 30°と270°を、ラジアンで表しましょう。（2）\nobreak $\dfrac{3\pi}{4}$ rad と 0.5\,rad を、度で表しましょう（0.5\,rad は小数第1位まで）。
\bkjans{（1）\nobreak $\dfrac{30\pi}{180}=\dfrac{\pi}{6}$ rad、$\dfrac{270\pi}{180}=\dfrac{3\pi}{2}$ rad。（2）\nobreak $\dfrac{3\times180^{\circ}}{4}=135^{\circ}$、$0.5\times\dfrac{180^{\circ}}{\pi}=\dfrac{90^{\circ}}{\pi}\approx28.6^{\circ}$。}
\end{jibun}
\bkivcpoint{換算でかける分数は、「消したい単位を分母に」。答えが出たら、目もりの図で、だいたい合っているか確かめましょう。1\,rad は約57°なので、2\,rad は120°より少し小さいはずです。}

\begin{reidai}[ラジアンで、弧の長さと面積を求める]
（1）\nobreak 半径\qty{5}{m}、中心角0.8\,rad のおうぎ形の、弧の長さと面積を求めましょう。\par
（2）\nobreak 弧の長さが\qty{6}{cm}、面積が\qty{15}{cm^2}のおうぎ形の、半径と中心角（rad）を求めましょう。\par
\kangaekata （1）$s=r\theta$、$S=\dfrac{1}{2}r^2\theta$ に入れます。（2）$S=\dfrac{1}{2}rs$ から $r$ を、$\theta=\dfrac{s}{r}$ から角を求めます。\par
\begin{bkcenter}
\begin{tikzpicture}[x=1mm, y=1mm]
  % （1）半径 5 m・0.8\,rad（約45.8°）、（2）半径 5 cm・1.2\,rad（約68.8°）。半径を 25 mm でかく
  \bkivcsector{(0,0)}{25}{0}{45.84}{bkOrange}
  \draw[bkInk, line width=0.5pt] (6,0) arc (0:45.84:6);
  % 第14回 B3：ラベルは、おうぎ形の線から 1.5 mm 以上はなす
  \node[font=\sffamily\footnotesize] at (18:16) {0.8 rad};
  \node[font=\sffamily\footnotesize, anchor=north] at (12.5,-1.4) {\qty{5}{m}};
  \node[font=\sffamily\footnotesize, anchor=west] at (22.9:27.5) {弧？};
  \node[font=\sffamily\small\bfseries, anchor=east] at (-2,20) {（1）};
  \begin{scope}[shift={(66,0)}]
    \bkivcsector{(0,0)}{25}{0}{68.75}{bkBlue}
    \draw[bkInk, line width=0.5pt] (6,0) arc (0:68.75:6);
    \node[font=\sffamily\footnotesize] at (20:10) {？};
    \node[font=\sffamily\footnotesize, anchor=north] at (12.5,-1.4) {半径？};
    \node[font=\sffamily\footnotesize, anchor=west] at (34.4:27.5) {弧 \qty{6}{cm}};
    \node[font=\sffamily\footnotesize] at (40:17) {\qty{15}{cm^2}};
    \node[font=\sffamily\small\bfseries, anchor=east] at (-2,20) {（2）};
  \end{scope}
\end{tikzpicture}
\end{bkcenter}

\shiki （1）$s=5\times0.8=4$\qquad $S=\dfrac{1}{2}\times5^2\times0.8=\dfrac{1}{2}\times25\times0.8=10$\par
\hspace*{2.6\zw}（2）$15=\dfrac{1}{2}\times r\times6$ より $15=3r$、$r=5$\par
\hspace*{4.3\zw}$\theta=\dfrac{s}{r}=\dfrac{6}{5}=1.2$\par
\hspace*{4.3\zw}確かめ：$\dfrac{1}{2}\times5^2\times1.2=\dfrac{1}{2}\times25\times1.2=15$\par
\kotae （1）\nobreak 弧の長さ\qty{4}{m}、面積\qty{10}{m^2}\quad（2）\nobreak 半径\qty{5}{cm}、中心角1.2\,rad
\end{reidai}

\begin{jibun}
（1）\nobreak 半径\qty{4}{cm}、中心角1.5\,rad のおうぎ形の、弧の長さと面積を求めましょう。（2）\nobreak 弧の長さが\qty{8}{cm}、面積が\qty{24}{cm^2}のおうぎ形の、半径と中心角（rad）を求めましょう。
\bkjans{（1）\nobreak $s=\qty{6}{cm}$、$S=\qty{12}{cm^2}$。（2）\nobreak $24=4r$ より $r=\qty{6}{cm}$。$\theta=\dfrac{4}{3}$ rad（約1.33\,rad）。}
\end{jibun}
\bkivcpoint{ラジアンなら、弧の長さは「半径×角」、面積は「半径×半径×角÷2」です。わからない量が2つあるときは、まず1つだけ求められる式を選びます（面積と弧から半径、次に半径と弧から角）。}

\begin{reidai}[動いた長さから、回った角を求める]
\noindent\begin{minipage}[t]{0.62\linewidth}
半径\qty{5}{cm}の小さいプーリー（ベルトをかける車）と、半径\qty{15}{cm}の大きいプーリーを、ベルトでつなぎます。ベルトが\qty{30}{cm}動いたとき、それぞれのプーリーは何 rad 回りますか。大きいプーリーの角は、度でも答えましょう。
\end{minipage}\hfill
\begin{minipage}[t]{0.34\linewidth}\vspace{0pt}\centering
\begin{tikzpicture}[x=0.85mm, y=0.85mm]
  % 小さいプーリー（半径 4 mm）と大きいプーリー（半径 12 mm）。外側の接線がベルト
  \pgfmathsetmacro\d{30}\pgfmathsetmacro\al{asin((12-4)/\d)}
  \draw[bkInk!70, line width=0.8pt, fill=bkGray!15] (0,0) circle (4);
  \draw[bkInk!70, line width=0.8pt, fill=bkGray!15] (\d,0) circle (12);
  \draw[bkOrange!80!black, line width=1.2pt] ({90+\al}:4) -- ($(\d,0)+({90+\al}:12)$);
  \draw[bkOrange!80!black, line width=1.2pt] ({-90-\al}:4) -- ($(\d,0)+({-90-\al}:12)$);
  \fill[bkInk] (0,0) circle (0.7); \fill[bkInk] (\d,0) circle (0.9);
  \node[font=\sffamily\footnotesize, anchor=north] at (0,-6.5) {\qty{5}{cm}};
  \node[font=\sffamily\footnotesize, anchor=north] at (\d,-13.5) {\qty{15}{cm}};
\end{tikzpicture}
\end{minipage}\par
\kangaekata ベルトがすべらなければ、どちらのプーリーのふちも、ベルトと同じ\qty{30}{cm}だけ動きます。動いた長さが弧の長さ $s$ なので、$\theta=\dfrac{s}{r}$ です。\par
\shiki 小さいプーリー：$\theta=\dfrac{30}{5}=6$\par
\vspace{2pt}\hspace*{2.6\zw}大きいプーリー：$\theta=\dfrac{30}{15}=2$\par\vspace{2pt}
\hspace*{2.6\zw}$2\ \text{rad}\times\dfrac{180^{\circ}}{\pi\ \text{rad}}=\dfrac{360^{\circ}}{\pi}\approx114.6^{\circ}$\par
\kotae 小さいプーリーは6\,rad、大きいプーリーは2\,rad（約114.6°）
\end{reidai}

\begin{jibun}
半径\qty{4}{cm}のプーリーと、半径\qty{12}{cm}のプーリーをベルトでつなぎます。ベルトが\qty{24}{cm}動いたとき、それぞれのプーリーは何 rad 回りますか。大きいプーリーの角は、度でも答えましょう。
\bkjans{小さいほうは $\dfrac{24}{4}=6$ rad、大きいほうは $\dfrac{24}{12}=2$ rad（約114.6°）。}
\end{jibun}
\bkivcpoint{ベルトでつながった車は、ふちが動く長さ $s$ が同じです。だから、半径が大きいほど、回る角は小さくなります（半径と角は反比例。2-15）。}

\begin{reidai}[タイヤが進む道のり]
半径\qty{0.33}{m}の自転車のタイヤが、すべらずに転がります。（1）\nobreak 100回転すると、何m進みますか（小数第1位まで）。（2）\nobreak \qty{1000}{m}進むには、約何回転しますか（整数で）。（3）\nobreak タイヤが1秒に3回転するとき、自転車の速さは時速約何kmですか（小数第1位まで）。\par
\kangaekata すべらずに転がるタイヤは、1回転で、地面に円周1つ分をおしつけながら進みます。1回転は $2\pi$ rad なので、$s=r\theta$ で $\theta=2\pi$ と考えても同じです。\par
\begin{bkcenter}
\begin{tikzpicture}[x=1mm, y=1mm]
  % 1 m を 25 mm でかく。半径 0.33 m → 8.25 mm、1回転 2.07 m → 51.8 mm
  \draw[bkGray, line width=1pt] (-12,0) -- (68,0);
  \foreach \xx/\op in {0/100, 51.8/45} {
    \draw[bkInk!80, line width=2.2pt, opacity=\op/100] (\xx,8.25) circle (7.6);
    \fill[bkInk!60, opacity=\op/100] (\xx,8.25) circle (1.2);
    \draw[bkRed, line width=0.9pt, opacity=\op/100] (\xx,8.25) -- (\xx,0.9);
    \fill[bkRed] (\xx,0) circle (1);}
  \draw[bkOrange, line width=1.2pt, -{Stealth[length=3pt]}] (0,-3) -- (51.8,-3);
  \node[font=\sffamily\footnotesize, text=bkOrange!80!black, anchor=north] at (25.9,-4.2) {1回転で $2\pi\times0.33\approx\qty{2.07}{m}$（赤い点が、ふたたび地面にふれる）};
\end{tikzpicture}
\end{bkcenter}

\shiki 1回転で進む道のり：$s=0.33\times2\pi=0.66\pi\approx2.0735$（m）\par
\hspace*{2.6\zw}（1）100回転：$0.66\pi\times100=66\pi\approx207.3$（m）\par
\hspace*{2.6\zw}（2）\qty{1000}{m}では：$1000\div2.0735\approx482.3$（回転）\par
\hspace*{2.6\zw}（3）1秒に3回転：$2.0735\times3\approx6.22$ で、\qty{6.22}{m/s}\par
\hspace*{4.3\zw}\qty{1}{m/s}＝\qty{3.6}{km/h}（2-10）なので、単位つきの1をかける\par
\hspace*{4.3\zw}$\qty{6.22}{m/s}\times\dfrac{\qty{3.6}{km/h}}{\qty{1}{m/s}}\approx\qty{22.4}{km/h}$\par
\kotae （1）\nobreak 約\qty{207.3}{m}\quad（2）\nobreak 約482回転\quad（3）\nobreak 時速約\qty{22.4}{km}
\end{reidai}

\begin{jibun}
半径\qty{0.3}{m}のタイヤが、すべらずに転がります。（1）\nobreak 50回転すると、何m進みますか（小数第1位まで）。（2）\nobreak \qty{600}{m}進むには、約何回転しますか（整数で）。（3）\nobreak 1秒に4回転するとき、速さは時速約何kmですか（小数第1位まで）。
\bkjans{1回転で $0.6\pi\approx\qty{1.885}{m}$。（1）\nobreak $30\pi\approx\qty{94.2}{m}$。（2）\nobreak $600\div1.885\approx318.3$ で、約318回転。（3）\nobreak $1.885\times4\approx\qty{7.54}{m/s}$。$7.54\times3.6\approx27.1$ で、時速約\qty{27.1}{km}。}
\end{jibun}
\bkivcpoint{1回転は $2\pi$ rad です。「何回転するか」は、「進む道のり÷1回転で進む道のり」で求めます。}

\begin{reidai}[接線と半径の角]
\noindent\begin{minipage}[t]{0.6\linewidth}
円Oの外の点Pから、円に2本の接線を引き、接点をA、Bとします。$\angle\mathrm{APB}$＝50°のとき、$\angle\mathrm{AOB}$を求めましょう。\par
\kangaekata 接線は接点で半径に垂直なので、$\angle\mathrm{OAP}$＝$\angle\mathrm{OBP}$＝90°です。四角形OAPBの内角の和は360°です（4-1：$n$ 角形の内角の和は $180^{\circ}\times(n-2)$）。
\end{minipage}\hfill
\begin{minipage}[t]{0.36\linewidth}\vspace{0pt}\centering
\begin{tikzpicture}[x=1mm, y=1mm]
  % 円 O(0,0) 半径 12。P は (12/sin25°, 0) ≈ (28.4, 0)
  \pgfmathsetmacro\pp{12/sin(25)}
  \draw[bkInk!60, line width=0.8pt] (0,0) circle (12);
  \draw[bkOrange, line width=1.1pt] (\pp,0) -- (65:12);
  \draw[bkOrange, line width=1.1pt] (\pp,0) -- (-65:12);
  \draw[bkBlue, line width=0.9pt] (65:12) -- (0,0) -- (-65:12);
  \bkivcright{(65:12)}{245}{335}
  \bkivcright{(-65:12)}{115}{25}
  \fill[bkInk] (0,0) circle (0.8);
  \node[font=\sffamily\footnotesize, anchor=east] at (-1,0) {O};
  \node[font=\sffamily\footnotesize, anchor=south] at (65:12.8) {A};
  \node[font=\sffamily\footnotesize, anchor=north] at (-65:12.8) {B};
  \node[font=\sffamily\footnotesize, anchor=west] at (\pp+0.8,0) {P};
  \node[font=\sffamily\footnotesize] at (\pp-6.5,0) {50°};
\end{tikzpicture}
\end{minipage}\par
\shiki 四角形OAPBの内角の和：$180^{\circ}\times(4-2)=360^{\circ}$\par
\hspace*{2.6\zw}$\angle\mathrm{AOB}=360^{\circ}-90^{\circ}-90^{\circ}-50^{\circ}=130^{\circ}$\par
\hspace*{2.6\zw}確かめ：$130^{\circ}+90^{\circ}+90^{\circ}+50^{\circ}=360^{\circ}$\par
\hspace*{2.6\zw}直角が2つあるので、$\angle\mathrm{AOB}$ と $\angle\mathrm{APB}$ の和は、いつも $360^{\circ}-180^{\circ}=180^{\circ}$ になる\par
\kotae 130°
\end{reidai}

\begin{jibun}
円Oの外の点Pから2本の接線を引き、接点をA、Bとします。$\angle\mathrm{APB}$＝70°のとき、$\angle\mathrm{AOB}$を求めましょう。
\bkjans{$\angle\mathrm{OAP}$＝$\angle\mathrm{OBP}$＝90°。$360^{\circ}-90^{\circ}-90^{\circ}-70^{\circ}=110^{\circ}$。}
\end{jibun}
\bkivcpoint{接線が出てきたら、まず接点と中心を結ぶ半径をかき、直角の印をつけましょう。それが、接線の問題を解く第一歩です。}

\begin{reidai}[円周角を求める]
（1）\nobreak 円Oで、弧ABに対する中心角$\angle\mathrm{AOB}$が100°です。弧ABに対する円周角$\angle\mathrm{APB}$を求めましょう。\par
（2）\nobreak ABが円の直径で、Pは円の上の点です。$\angle\mathrm{PAB}$＝35°のとき、$\angle\mathrm{PBA}$を求めましょう。\par
\kangaekata （1）円周角は中心角の半分。（2）直径に対する円周角は90°なので、$\angle\mathrm{APB}$＝90°。三角形PABの内角の和は180°です。\par
\begin{bkcenter}
\begin{tikzpicture}[x=1mm, y=1mm]
  % （1）A は 220°、B は 320°（中心角 100°）、P は 90°
  \draw[bkInk!60, line width=0.8pt] (0,0) circle (15);
  \draw[bkInk, line width=0.7pt] (220:15) -- (0,0) -- (320:15);
  \draw[bkBlue, line width=0.8pt] (220:15) -- (90:15) -- (320:15);
  \draw[bkInk, line width=0.5pt] (220:4) arc (220:320:4);
  \node[font=\sffamily\footnotesize] at (270:7.2) {100°};
  \fill[bkInk] (0,0) circle (0.7);
  \foreach \a/\l/\an in {220/A/north east, 320/B/north west, 90/P/south} {\fill[bkInk] (\a:15) circle (0.8); \node[font=\sffamily\footnotesize, anchor=\an] at (\a:15.6) {\l};}
  \node[font=\sffamily\footnotesize, text=bkBlue!80!black] at (0,8) {？};
  \node[font=\sffamily\small\bfseries, anchor=east] at (-17,12) {（1）};
  % （2）AB は直径、P は 70°（∠PAB＝35°）
  \begin{scope}[shift={(62,0)}]
    \draw[bkInk!60, line width=0.8pt] (0,0) circle (15);
    \draw[bkOrange, line width=1pt] (-15,0) -- (15,0);
    \draw[bkBlue, line width=0.8pt] (-15,0) -- (70:15) -- (15,0);
    \fill[bkInk] (0,0) circle (0.7);
    \draw[bkInk, line width=0.5pt] (-15,0) ++(0:6) arc (0:35:6);
    \node[font=\sffamily\footnotesize] at ($(-15,0)+(17.5:9.2)$) {35°};
    \node[font=\sffamily\footnotesize, text=bkBlue!80!black] at ($(15,0)+(152:6.5)$) {？};
    \foreach \p/\l/\an in {{(-15,0)}/A/east, {(15,0)}/B/west, {(70:15)}/P/south} {\node[font=\sffamily\footnotesize, anchor=\an] at \p {\l};}
    \node[font=\sffamily\small\bfseries, anchor=east] at (-17,12) {（2）};
  \end{scope}
\end{tikzpicture}
\end{bkcenter}

\shiki （1）$100^{\circ}\div2=50^{\circ}$\par
\hspace*{2.6\zw}（2）$\angle\mathrm{APB}$＝90°。$\angle\mathrm{PBA}=180^{\circ}-90^{\circ}-35^{\circ}=55^{\circ}$\par
\kotae （1）\nobreak 50°\quad（2）\nobreak 55°
\end{reidai}

\begin{jibun}
（1）\nobreak 弧ABに対する中心角が140°のとき、円周角を求めましょう。（2）\nobreak ABが円の直径で、Pは円の上の点です。$\angle\mathrm{PAB}$＝28°のとき、$\angle\mathrm{PBA}$を求めましょう。
\bkjans{（1）\nobreak $140^{\circ}\div2=70^{\circ}$。（2）\nobreak $\angle\mathrm{APB}$＝90°。$180^{\circ}-90^{\circ}-28^{\circ}=62^{\circ}$。}
\end{jibun}
\bkivcpoint{円周角の問題では、① 同じ弧に対する中心角をさがす、② 直径があれば90°の角をさがす。そのあとで、三角形の内角の和180°を使います。}

\begin{reidai}[3点を通る円の中心（外心）を求める]
\noindent\begin{minipage}[t]{0.6\linewidth}
座標平面の3点 A$(0,\ 0)$、B$(6,\ 0)$、C$(0,\ 4)$ を通る円の、中心の座標を求めましょう。\par
\kangaekata 中心は、辺ABの垂直二等分線と、辺ACの垂直二等分線の交わる点です。ABは $x$ 軸の上にあるので、垂直二等分線は、ABのまん中を通る、たての直線です。ACは $y$ 軸の上にあるので、垂直二等分線は、ACのまん中を通る、横の直線です。
\end{minipage}\hfill
\begin{minipage}[t]{0.36\linewidth}\vspace{0pt}\centering
\begin{tikzpicture}
\begin{axis}[bkplot, width=46mm, height=40mm, xmin=-1, xmax=7.4, ymin=-1, ymax=5.4, xtick={3,6}, ytick={2,4}, xlabel={$x$}, ylabel={$y$}]
  \draw[bkGray!70, line width=0.6pt] (axis cs:3,2) circle [radius={sqrt(13)}];
  \draw[bkInk, line width=0.9pt] (axis cs:0,0) -- (axis cs:6,0) -- (axis cs:0,4) -- cycle;
  \draw[bkOrange, line width=0.8pt, dash pattern=on 2pt off 1pt] (axis cs:3,-0.8) -- (axis cs:3,4.6);
  \draw[bkPurple, line width=0.8pt, dash pattern=on 2pt off 1pt] (axis cs:-0.8,2) -- (axis cs:6.6,2);
  \addplot[only marks, mark=*, mark size=1.5pt, bkInk] coordinates {(0,0) (6,0) (0,4) (3,2)};
  \node[font=\sffamily\footnotesize, anchor=north east] at (axis cs:0,0) {A};
  \node[font=\sffamily\footnotesize, anchor=south west] at (axis cs:6,0) {B};
  \node[font=\sffamily\footnotesize, anchor=south west] at (axis cs:0,4) {C};
\end{axis}
\end{tikzpicture}
\end{minipage}\par
\shiki ABのまん中は $(3,\ 0)$。垂直二等分線は、$x=3$ のたての直線\par
\hspace*{2.6\zw}ACのまん中は $(0,\ 2)$。垂直二等分線は、$y=2$ の横の直線\par
\hspace*{2.6\zw}2本が交わる点は $(3,\ 2)$\par
\hspace*{2.6\zw}確かめ：$(3,\ 2)$ から見ると、\par
\hspace*{4.3\zw}Aは、左へ3・下へ2の所\par
\hspace*{4.3\zw}Bは、右へ3・下へ2の所\par
\hspace*{4.3\zw}Cは、左へ3・上へ2の所\par
\hspace*{4.3\zw}どれも「横に3・たてに2」の動きなので、3つの点までの長さは等しい\par
\hspace*{4.3\zw}（長さそのものは、4-4で計算できる）\par
\hspace*{2.6\zw}確かめ：$\angle\mathrm{CAB}$＝90°なので、BCは直径のはず。BCのまん中は $\left(\dfrac{6+0}{2},\ \dfrac{0+4}{2}\right)=(3,\ 2)$ で合う\par
\kotae $(3,\ 2)$
\end{reidai}

\begin{jibun}
3点 A$(0,\ 0)$、B$(8,\ 0)$、C$(0,\ 6)$ を通る円の、中心の座標を求めましょう。
\bkjans{ABの垂直二等分線は $x=4$、ACの垂直二等分線は $y=3$。中心は $(4,\ 3)$（BCのまん中と同じ）。}
\end{jibun}
\bkivcpoint{外心は「3つの頂点から同じ長さの点」です。垂直二等分線を2本引けば決まります（3本めも同じ点を通ります）。直角三角形なら、外心は斜辺のまん中です。}

\begin{reidai}[極座標で、点の場所を言う・点をかく]
\noindent\begin{minipage}[t]{0.5\linewidth}
右の図は、中心Oのまわりに、中心からの距離が1・2・3・4の円と、30°ごとの線をかいたものです。\par
（1）\nobreak 点A、B、Cの極座標を答えましょう（角は度で）。\par
（2）\nobreak 極座標が $\left(2,\ \dfrac{\pi}{2}\right)$ の点Pと、$\left(3,\ \dfrac{7\pi}{6}\right)$ の点Qは、図のどこにありますか。
\end{minipage}\hfill
\begin{minipage}[t]{0.46\linewidth}\vspace{0pt}\centering
\begin{tikzpicture}[x=1mm, y=1mm]
  \bkivcpolar{5}{4}
  \foreach \r/\a/\n/\dx/\dy/\an in {3/60/A/1.8/0/west, 4/150/B/0/-1.8/north, 2/300/C/1.8/0/west} {
    \fill[bkBlue!80!black] (\a:{5*\r}) circle (1.1);
    \node[font=\sffamily\footnotesize\bfseries, text=bkBlue!80!black, fill=white, inner sep=0.6pt, anchor=\an] at ($(\a:{5*\r})+(\dx,\dy)$) {\n};}
\end{tikzpicture}
\end{minipage}\par
\kangaekata $r$ は「中心から何番めの円の上か」、$\theta$ は「0°の線から左回りに、何度の線の上か」で読みます。（2）は、ラジアンを度に直してから（④）、その向きへ $r$ だけ進んだ点をさがします。\par
\shiki （1）A：3番めの円と、60°の線が交わる点。$(3,\ 60°)$\par
\hspace*{4.3\zw}B：4番めの円と、150°の線が交わる点。$(4,\ 150°)$\par
\hspace*{4.3\zw}C：2番めの円と、300°の線が交わる点。$(2,\ 300°)$\par
\hspace*{2.6\zw}（2）P：$\dfrac{\pi}{2}\ \text{rad}\times\dfrac{180^{\circ}}{\pi\ \text{rad}}=90^{\circ}$。真上の向きへ2\par
\hspace*{4.3\zw}Q：$\dfrac{7\pi}{6}\ \text{rad}\times\dfrac{180^{\circ}}{\pi\ \text{rad}}=\dfrac{7\times180^{\circ}}{6}=210^{\circ}$\par
\hspace*{4.3\zw}210°は、真左（180°）から、さらに30°左回りの向き\par
\hspace*{4.3\zw}その向き（左下）へ3\par
\kotae （1）\nobreak A $(3,\ 60°)$、B $(4,\ 150°)$、C $(2,\ 300°)$\par
\hspace*{2.6\zw}（2）\nobreak P：2番めの円と90°の線が交わる点（Oの真上）\par
\hspace*{4.3\zw}Q：3番めの円と210°の線が交わる点（左下）
\end{reidai}

\begin{jibun}
\noindent\begin{minipage}[t]{0.56\linewidth}
（1）\nobreak 右の図の点D、E、Fの極座標を答えましょう（円は中心から1・2・3・4、線は30°ごと）。（2）\nobreak 極座標が $(1,\ \pi)$ の点Rと、$\left(3,\ \dfrac{5\pi}{3}\right)$ の点Sは、図のどこにありますか。
\end{minipage}\hfill
\begin{minipage}[t]{0.4\linewidth}\vspace{0pt}\centering
\begin{tikzpicture}[x=0.85mm, y=0.85mm]
  \bkivcpolar{5}{4}
  \foreach \r/\a/\n/\dx/\dy/\an in {2/30/D/1.8/0/west, 3/120/E/1.8/0/west, 4/240/F/1.8/0/west} {
    \fill[bkBlue!80!black] (\a:{5*\r}) circle (1.1);
    \node[font=\sffamily\footnotesize\bfseries, text=bkBlue!80!black, fill=white, inner sep=0.6pt, anchor=\an] at ($(\a:{5*\r})+(\dx,\dy)$) {\n};}
\end{tikzpicture}
\end{minipage}
\bkjans{（1）\nobreak D $(2,\ 30°)$、E $(3,\ 120°)$、F $(4,\ 240°)$。（2）\nobreak $\pi$ は180°なので、Rは1番めの円と180°の線が交わる点（Oの真左）。$\dfrac{5\pi}{3}$ は300°なので、Sは3番めの円と300°の線が交わる点（右下）。}
\end{jibun}
\bkivcpoint{極座標は「距離、角」の順。角は0°の線（右）から左回りに読みます。ラジアンで書いてあったら、度に直すと、どの線の上かがすぐわかります。}

\begin{motto}[2]
\raggedright
\textbf{地図や船の「方位角」は、北から右回り}\par
地図や船・飛行機では、向きを「北から右回りに何度」で言うことが多く、これを\textbf{方位角}（ほういかく）といいます。東は方位角90°、南は180°です。この本の極座標は「東（右）から左回り」なので、同じ東の向きが0°、北が90°になります。北東だけは、どちらで言っても45°です。角の数を見たら、「どこから、どちら回りにはかった角か」を、まず確かめましょう。
\par\vspace{3pt}
{\centering
\begin{tikzpicture}[x=1mm, y=1mm]
  \foreach \xs/\ttl in {0/この本の極座標, 62/地図の方位角} {
    \begin{scope}[xshift=\xs mm]
      \draw[bkGray!70, line width=0.5pt] (0,0) circle (11);
      \draw[bkGray, line width=0.5pt] (-11,0) -- (11,0) (0,-11) -- (0,11);
      \node[font=\sffamily\footnotesize, anchor=south] at (0,12.6) {北};
      \node[font=\sffamily\footnotesize, anchor=west] at (12.6,0) {東};
      \node[font=\sffamily\footnotesize, anchor=north] at (0,-12.6) {南};
      \node[font=\sffamily\footnotesize, anchor=east] at (-12.6,0) {西};
      \node[font=\sffamily\bfseries\footnotesize, anchor=north] at (0,-18.5) {\ttl};
      \fill[bkInk] (0,0) circle (0.7);
    \end{scope}}
  % この本の極座標：東が 0°、左回り
  \draw[bkOrange, line width=1.1pt, -{Stealth[length=3.5pt]}] (6,0) arc (0:85:6);
  \node[font=\sffamily\footnotesize, text=bkOrange!80!black, anchor=west] at (25,7) {東が0°};
  \node[font=\sffamily\footnotesize, text=bkOrange!80!black, anchor=west] at (25,1.5) {左回り};
  \node[font=\sffamily\footnotesize, text=bkOrange!80!black, anchor=west] at (25,-4) {北は90°};
  % 方位角：北が 0°、右回り
  \draw[bkBlue, line width=1.1pt, -{Stealth[length=3.5pt]}] (62,6) arc (90:5:6);
  \node[font=\sffamily\footnotesize, text=bkBlue!80!black, anchor=west] at (87,7) {北が0°};
  \node[font=\sffamily\footnotesize, text=bkBlue!80!black, anchor=west] at (87,1.5) {右回り};
  \node[font=\sffamily\footnotesize, text=bkBlue!80!black, anchor=west] at (87,-4) {東は90°};
\end{tikzpicture}\par}
\end{motto}

%% ---------------- まちがえやすいところ ----------------
\Needspace{14\baselineskip}
\bksection{tsumazuki}{まちがえやすいところをチェック}

\begin{tsumazuki}{換算でかける分数が、さかさま}
\bkng{$60^{\circ}=60\times\dfrac{180}{\pi}$ rad}
\bkok{$60^{\circ}=60^{\circ}\times\dfrac{\pi\ \text{rad}}{180^{\circ}}=\dfrac{\pi}{3}$ rad}
単位をつけて書けば、°が分母にあるほうをかけると°が消えることがわかります。×の式は約3438になり、1周の $2\pi\approx6.28$ より大きいので、おかしいと気づけます。
\end{tsumazuki}

\begin{tsumazuki}{$s=r\theta$ に、度のまま入れる}
\bkng{半径2、中心角30°の弧の長さは $2\times30=60$}
\bkok{$30^{\circ}=\dfrac{\pi}{6}$ rad なので、$s=2\times\dfrac{\pi}{6}=\dfrac{\pi}{3}\approx1.05$}
$s=r\theta$ は、$\theta$ がラジアンのときの式です。円周は $2\pi\times2\approx12.6$ なのに、×の60は円周より長くなっています。
\end{tsumazuki}

\begin{tsumazuki}{半径と直径を取りちがえる}
\bkng{半径5の円の円周は $5\pi$}
\bkok{円周＝$2\pi r=2\pi\times5=10\pi$}
2-9の「直径×$\pi$」の直径は、半径の2倍です。$r$ は半径なので、2をかけるのをわすれないこと。
\end{tsumazuki}

\begin{tsumazuki}{円周角を、中心角と同じや2倍にする}
\bkng{中心角80°なので、円周角は80°（または160°）}
\bkok{円周角は中心角の半分。$80^{\circ}\div2=40^{\circ}$}
円の上の点Pは、中心Oよりも弧ABから遠いので、ABを見こむ角は小さくなります。分度器ではかって確かめてみよう。
\end{tsumazuki}

\begin{tsumazuki}{極座標の角を、真上からはかる}
\bkng{Oの真上の向きへ3の点は $(3,\ 0°)$}
\bkok{真上は、右（0°）から左回りに90°なので $(3,\ 90°)$}
極座標の角は、$x$ 軸の正の向き（右）を0°にして、左回りにはかります。地図の方位角（北から右回り）とまちがえないようにしましょう。
\end{tsumazuki}

\begin{tsumazuki}{1\,rad を、1°くらいの小さい角だと思う}
\bkng{1\,rad は、1°とほとんど同じ大きさ}
\bkok{1\,rad は約57.3°。1周は約6.28\,rad}
1\,rad は、弧の長さが半径と同じになる角です。半径を6本ちょっとならべると1周するので、1\,rad は1周の6分の1より少し小さい角です。
\end{tsumazuki}

%% ---------------- 物理・工学 ----------------
\begin{butsuri}
\par\addvspace{3pt}\textbf{タイヤとスピードメーター}\par\noindent
\begin{minipage}[c]{0.68\linewidth}\raggedright
自動車のスピードメーターは、タイヤ（車輪）が1秒に何回転したかを数えて、速さを計算しています。半径\qty{0.30}{m}のタイヤなら、1回転で約\qty{1.885}{m}（$2\pi\times0.30$）進みます。1秒に10回転していれば、1秒に\qty{18.85}{m}進みます。時速にすると、単位つきの1（2-10）をかけて、
\[ \qty{18.85}{m/s}\times\frac{\qty{3.6}{km/h}}{\qty{1}{m/s}}\approx\qty{67.9}{km/h} \]
です。タイヤがすりへって半径が小さくなると、1回転で進む道のりが短くなるので、メーターは実際より少し速い速さを表示します。
\end{minipage}\hfill
\begin{minipage}[c]{0.28\linewidth}\centering
\begin{tikzpicture}[x=0.85mm, y=0.85mm]
  \draw[bkGray, line width=1pt] (-18,0) -- (20,0);
  \draw[bkInk!80, line width=2.6pt] (0,12) circle (10.5);
  \draw[bkInk!40, line width=0.7pt] (0,12) circle (6);
  \fill[bkInk!60] (0,12) circle (1.6);
  \draw[bkOrange, line width=1.2pt] (0,12) -- (12,12);
  \node[font=\sffamily\footnotesize, anchor=south] at (6,12.4) {$r$};
  \draw[bkBlue, line width=1.1pt, -{Stealth[length=3pt]}] (-12,-3) -- (14,-3);
  \node[font=\sffamily\footnotesize, text=bkBlue!80!black, anchor=north] at (1,-4) {1回転で $2\pi r$};
\end{tikzpicture}
\end{minipage}

\par\addvspace{4pt}\textbf{ベルトとプーリー}\par\noindent
\begin{minipage}[c]{0.6\linewidth}\raggedright
エンジンや洗たく機のモーターは、ベルトとプーリーで回転を伝えます。ベルトがすべらなければ、2つのプーリーのふちは同じ長さだけ動きます。$s=r\theta$ で $s$ が同じなので、半径3倍のプーリーは、3分の1の角しか回りません（例題4）。半径のちがうプーリーを組み合わせると、回転をおそくして、そのかわりに回す力を大きくできます。

たて型の洗たく機では、モーターの小さいプーリーと、洗たく槽の大きいプーリー（半径が約6倍）をベルトでつなぎ、回転を約 $\dfrac{1}{6}$ におとしています。
\end{minipage}\hfill
\begin{minipage}[c]{0.36\linewidth}\centering
\begin{tikzpicture}[x=0.8mm, y=0.8mm]
  % 小さいプーリー（半径 3）と大きいプーリー（半径 15）。外側の接線がベルト
  \pgfmathsetmacro\d{30}\pgfmathsetmacro\al{asin((15-3)/\d)}
  \draw[bkInk!70, line width=0.8pt, fill=bkGray!15] (0,0) circle (3);
  \draw[bkInk!70, line width=0.8pt, fill=bkGray!15] (\d,0) circle (15);
  \draw[bkOrange!80!black, line width=1.2pt] ({90+\al}:3) -- ($(\d,0)+({90+\al}:15)$);
  \draw[bkOrange!80!black, line width=1.2pt] ({-90-\al}:3) -- ($(\d,0)+({-90-\al}:15)$);
  \fill[bkInk] (0,0) circle (0.7); \fill[bkInk] (\d,0) circle (0.9);
  \node[font=\sffamily\footnotesize, anchor=north] at (0,-5.5) {モーター};
  \node[font=\sffamily\footnotesize, anchor=north] at (\d,-16.5) {洗たく槽};
  \node[font=\sffamily\footnotesize, text=bkOrange!80!black, anchor=south] at (14,10.4) {ベルト};
\end{tikzpicture}
\end{minipage}
\end{butsuri}

\begin{butsuri}
\par\addvspace{3pt}\textbf{ハンマー投げと、回る傘のしずく}\par\noindent
\begin{minipage}[c]{0.62\linewidth}
ハンマー投げでは、ワイヤーの先のハンマーを、半径約\qty{2}{m}の円をえがくように回します。ワイヤーは、ハンマーを中心のほうへ引っぱっています。手をはなすと、ハンマーはその瞬間に動いていた向き、つまり\textbf{円の接線の向き}（ワイヤーに垂直）に飛んでいきます。ぬれた傘を回すと、しずくがふちから接線の向きに飛ぶのも、同じしくみです。
\end{minipage}\hfill
\begin{minipage}[c]{0.34\linewidth}\centering
\begin{tikzpicture}[x=1mm, y=1mm]
  \draw[bkInk!40, line width=0.6pt, dash pattern=on 2pt off 1.5pt] (0,0) circle (16);
  \fill[bkInk] (0,0) circle (1);
  \draw[bkInk!70, line width=0.9pt] (0,0) -- (40:16);
  \fill[bkGray!80!black] (40:16) circle (2.4);
  \draw[bk arrow, bkOrange, line width=1.5pt] ($(40:16)+(130:3)$) -- ($(40:16)+(130:17)$);
  \bkivcright{(40:16)}{220}{130}
  \node[font=\sffamily\footnotesize, text=bkOrange!80!black, anchor=south east] at ($(40:16)+(130:14)+(0,1)$) {接線の向き};
  \node[font=\sffamily\footnotesize, anchor=north west] at ($(20:8)+(0,-0.5)$) {ワイヤー};
\end{tikzpicture}
\end{minipage}
\end{butsuri}

%% ---------------- 章末問題 ----------------
\bkendproblems
\bkproblemgroup{基本}{2}

\bkproblem 半径\qty{9}{cm}、中心角40°のおうぎ形の、（1）\nobreak 弧の長さと、（2）\nobreak 面積を求めましょう。$\pi$ を使って表し、小数（約）でも答えましょう。（\bkreidairef{1}）
\bkpans{（1）\nobreak $18\pi\times\dfrac{40}{360}=2\pi$\,cm（約\qty{6.28}{cm}）。（2）\nobreak $\dfrac{1}{2}\times9\times2\pi=9\pi$\,\unit{cm^2}（約\qty{28.3}{cm^2}）。}

\bkproblem （1）\nobreak 120°と225°を、ラジアンで表しましょう。（2）\nobreak $\dfrac{5\pi}{3}$ rad と 3\,rad を、度で表しましょう（3\,rad は小数第1位まで）。（\bkreidairef{2}）
\bkpans{（1）\nobreak $\dfrac{2\pi}{3}$ rad、$\dfrac{5\pi}{4}$ rad。（2）\nobreak 300°、$\dfrac{540^{\circ}}{\pi}\approx171.9^{\circ}$。}

\bkproblem （1）\nobreak 半径\qty{10}{cm}、中心角 $\dfrac{\pi}{6}$ rad のおうぎ形の、弧の長さと面積を求めましょう（$\pi$ を使って）。（2）\nobreak 弧の長さが\qty{10}{cm}、面積が\qty{40}{cm^2}のおうぎ形の、半径と中心角を求めましょう。（\bkreidairef{3}）
\bkpans{（1）\nobreak $s=\dfrac{5\pi}{3}$\,cm（約\qty{5.24}{cm}）、$S=\dfrac{25\pi}{3}$\,\unit{cm^2}（約\qty{26.2}{cm^2}）。（2）\nobreak $40=5r$ より $r=\qty{8}{cm}$、$\theta=1.25$ rad。}

\bkproblem 半径\qty{6}{cm}と\qty{18}{cm}のプーリーをベルトでつなぎます。ベルトが\qty{36}{cm}動いたとき、それぞれ何 rad 回りますか。大きいほうは度でも答えましょう。（\bkreidairef{4}）
\bkpans{小さいほうは $\dfrac{36}{6}=6$ rad、大きいほうは $\dfrac{36}{18}=2$ rad（約114.6°）。}

\bkproblem 半径\qty{0.35}{m}のタイヤがすべらずに転がります（\bkreidairef{5}）。（1）\nobreak 200回転で何m進む？（小数第1位まで）（2）\nobreak \qty{1000}{m}では約何回転？（整数で）
\bkpans{1回転で $0.7\pi\approx\qty{2.199}{m}$。（1）\nobreak $140\pi\approx\qty{439.8}{m}$。（2）\nobreak $1000\div2.199\approx454.7$ で、約455回転。}

\bkproblem 次の角を求めましょう（\bkreidairef{6}・\bkreidairef{7}）。（1）\nobreak 円Oの外の点Pから引いた2本の接線の接点がA、Bで、$\angle\mathrm{APB}$＝40°のときの$\angle\mathrm{AOB}$。（2）\nobreak 中心角76°の弧に対する円周角。（3）\nobreak ABが直径、Pが円の上の点で、$\angle\mathrm{PBA}$＝64°のときの$\angle\mathrm{PAB}$。
\bkpans{（1）\nobreak $360^{\circ}-90^{\circ}-90^{\circ}-40^{\circ}=140^{\circ}$。（2）\nobreak $76^{\circ}\div2=38^{\circ}$。（3）\nobreak $\angle\mathrm{APB}$＝90°。$180^{\circ}-90^{\circ}-64^{\circ}=26^{\circ}$。}

\bkproblem 例題9と同じ図をノートにかき、極座標が $\left(3,\ \dfrac{2\pi}{3}\right)$ の点Pと $\left(2,\ \dfrac{3\pi}{2}\right)$ の点Qをかきましょう。角は何度ですか。\mbox{（\bkreidairef{9}）}
\bkpans{Pの角は $\dfrac{2\times180^{\circ}}{3}=120^{\circ}$ で、Pは3番めの円と120°の線が交わる点（左上）。Qの角は $\dfrac{3\times180^{\circ}}{2}=270^{\circ}$ で、Qは2番めの円と270°の線が交わる点（中心の真下）。}

\bkproblemgroup{チャレンジ}{3}

\bkproblem A$(0,\ 0)$、B$(10,\ 0)$、C$(0,\ 4)$ を通る円の中心を求めましょう。\mbox{（\bkreidairef{8}）}
\bkpans{ABの垂直二等分線は $x=5$、ACの垂直二等分線は $y=2$。中心は $(5,\ 2)$。}

\bkproblem 半径\qty{0.32}{m}のタイヤが1秒に8回転しています。（1）\nobreak 1秒に何m進みますか（小数第2位まで）。（2）\nobreak 時速約何kmですか（整数で）。（\bkreidairef{5}）
\bkpans{（1）\nobreak 1回転で $0.64\pi\approx\qty{2.011}{m}$。8回転で $5.12\pi\approx\qty{16.08}{m}$。（2）\nobreak $16.08\times3.6\approx57.9$ で、時速約\qty{58}{km}。}

\bkproblem 半径\qty{12}{cm}、弧の長さ $4\pi$\,cm のおうぎ形の、（1）\nobreak 中心角（ラジアンと度）、（2）\nobreak 面積（$\pi$ を使って）を求めましょう。（\bkreidairef{2}・\bkreidairef{3}）
\bkpans{（1）\nobreak $\theta=\dfrac{4\pi}{12}=\dfrac{\pi}{3}$ rad、60°。（2）\nobreak $\dfrac{1}{2}\times12\times4\pi=24\pi$\,\unit{cm^2}（約\qty{75.4}{cm^2}）。}

\begin{butsuri}
\par\addvspace{3pt}\textbf{気象レーダーと、雨雲の地図}\par\noindent
\begin{minipage}[c]{0.7\linewidth}\raggedright
気象レーダーは、アンテナを回しながら電波を出し、雨つぶではね返った電波をはかります。もどるまでの時間で距離 $r$ が、アンテナの向きで角 $\theta$ がわかります（⑩）。電波は1秒に約30万km進むので、\qty{60}{km}先の雨なら約0.0004秒で往復します。集めた $(r,\ \theta)$ を色でぬると、雨雲の地図になります。
\end{minipage}\hfill
\begin{minipage}[c]{0.27\linewidth}\centering
\begin{tikzpicture}[x=0.85mm, y=0.85mm]
  \foreach \k in {1,2,3} {\draw[bkGray!55, line width=0.4pt] (0,0) circle ({6*\k});}
  \foreach \a in {0,30,...,330} {\draw[bkGray!55, line width=0.4pt] (0,0) -- (\a:18);}
  \fill[bkBlue!35, opacity=0.8] (40:10) ellipse (4.5 and 3);
  \fill[bkBlue!60, opacity=0.8] (40:10.5) ellipse (2 and 1.4);
  \fill[bkBlue!30, opacity=0.8] (200:13) ellipse (3.5 and 2.2);
  \draw[bkOrange, line width=1.2pt] (0,0) -- (40:10);
  \fill[bkInk] (0,0) circle (0.9);
  \node[font=\sffamily\footnotesize, anchor=north, inner sep=1pt] at (0,-19.5) {レーダーの画面};
\end{tikzpicture}
\end{minipage}
\end{butsuri}

\begin{furikaeri}
\begin{checklist}
\item 弦・弧・おうぎ形がわかり、弧の長さとおうぎ形の面積を、中心角に比例させて求められる。\bkfurref{bk:4-2:hirei}
\item 角をラジアンで表し、度と換算して、$s=r\theta$ と $S=\dfrac{1}{2}r^2\theta$ を使える。\bkfurref{bk:4-2:rad}
\item 接線が半径に垂直なわけを説明でき、DEG・RAD と極座標 $(r,\ \theta)$ を使える。\bkfurref{bk:4-2:sessen}
\end{checklist}
\end{furikaeri}

\bkfukidashi{delta}{回るものの道のりは $s=r\theta$。次の章は「比」だよ。}
